\documentclass[12pt]{article}
\usepackage[T1]{fontenc}
\usepackage[utf8]{inputenc}
\usepackage[english]{babel}
\usepackage{lmodern}
\usepackage{microtype}
\usepackage{amsmath,amssymb,amsthm,mathtools}
\usepackage{mathrsfs}
\usepackage{enumitem}
\usepackage{geometry}
\geometry{margin=1.02in}
\usepackage{setspace}
\setstretch{1.045}
\usepackage{graphicx}
\usepackage[all,cmtip]{xy}
\usepackage{hyperref}
\hypersetup{colorlinks=true,linkcolor=blue,citecolor=blue,urlcolor=blue,hypertexnames=false}

\newcommand{\U}{\mathcal{U}}
\newcommand{\V}{\mathcal{V}}
\newcommand{\F}{\mathcal{F}}
\newcommand{\G}{\mathcal{G}}
\newcommand{\Hh}{\mathcal{H}}
\newcommand{\Oo}{\mathcal{O}}
\newcommand{\A}{\mathcal{A}}
\newcommand{\B}{\mathcal{B}}
\newcommand{\C}{\mathcal{C}}
\newcommand{\E}{\mathcal{E}}
\newcommand{\M}{\mathcal{M}}
\newcommand{\N}{\mathcal{N}}
\newcommand{\J}{\mathcal{J}}
\newcommand{\Lcal}{\mathcal{L}}
\newcommand{\Zz}{\mathbb{Z}}
\newcommand{\Nn}{\mathbb{N}}
\newcommand{\Qq}{\mathbb{Q}}
\newcommand{\Rr}{\mathbb{R}}
\newcommand{\Cc}{\mathbb{C}}
\newcommand{\Pp}{\mathbb{P}}
\newcommand{\Aa}{\mathbb{A}}
\newcommand{\Gm}{\mathbb{G}_m}
\newcommand{\LL}{\mathcal L}
\newcommand{\Primes}{\mathcal P}
\newcommand{\Spec}{\operatorname{Spec}}
\newcommand{\Hom}{\operatorname{Hom}}
\newcommand{\End}{\operatorname{End}}
\newcommand{\Isom}{\operatorname{Isom}}
\newcommand{\Tor}{\operatorname{Tor}}
\newcommand{\Ker}{\operatorname{Ker}}
\newcommand{\Coker}{\operatorname{Coker}}
\newcommand{\Ob}{\operatorname{Ob}}
\newcommand{\Ass}{\operatorname{Ass}}
\newcommand{\Supp}{\operatorname{Supp}}
\newcommand{\Pic}{\operatorname{Pic}}
\newcommand{\Et}{\operatorname{Et}}

\newcommand{\Sch}{\operatorname{Sch}}
\newcommand{\Sh}{\operatorname{Sh}}
\newcommand{\Ga}{\mathbb{G}_a}
\newcommand{\PP}{\mathbb{P}}
\newcommand{\an}{\mathrm{an}}
\newcommand{\zar}{\mathrm{zar}}
\newcommand{\cd}{\operatorname{cd}}
\newcommand{\rad}{\operatorname{rad}}
\newcommand{\codim}{\operatorname{codim}}
\newcommand{\Char}{\operatorname{char}}
\newcommand{\red}{\mathrm{red}}
\newcommand{\ab}{\mathrm{ab}}
\newcommand{\et}{\mathrm{\acute{e}t}}
\newcommand{\etale}{\ifmmode\textnormal{\'etale}\else\'etale\fi}
\newcommand{\fprod}[3]{{#1}\times_{#3}{#2}}
\newcommand{\ie}{i.e.\ }
\newcommand{\cf}{cf.\ }
\newcommand{\resp}{respectively}
\newcommand{\etc}{etc.\ }
\newcommand{\loccit}{loc. cit.}
\newcommand{\SGA}{SGA}
\newcommand{\EGA}{EGA}
\newcommand{\qedtext}{\hfill\(\square\)}
\DeclareMathOperator{\Quot}{Quot}
\DeclareMathOperator{\Alg}{Alg}
\DeclareMathOperator{\trdeg}{trdeg}
\newcommand{\Hhom}{\mathcal{H}\!om}
\newenvironment{statement}[1]{\par\medskip\noindent\textbf{#1.}\ }{\par\medskip}
\newenvironment{remarktext}[1]{\par\medskip\noindent\textbf{#1.}\ }{\par\medskip}
\newenvironment{prooftext}[1][Proof]{\par\noindent\emph{#1.}\ }{\par\medskip}
\setlist[enumerate]{leftmargin=2.25em}
\setlist[itemize]{leftmargin=2.25em}
\allowdisplaybreaks

\begin{document}

\begin{center}
{\Large \textbf{SGA 4 Translation Project -- translation segment}}\\[0.4em]
{\large \textbf{Proper Base Change through Cohomology with Proper Supports}}\\[0.25em]
{\normalsize Draft English LaTeX translation; compile-checked, not yet line-proofed against the scan.}
\end{center}

\bigskip

% ============================================================
\begin{center}
{\Large \textbf{SGA 4, Expos\'e XII -- continuation}}\\[0.35em]
{\large \textbf{The Proper Base Change Theorem}}\\[0.25em]
{\normalsize M. Artin}
\end{center}

\section*{6. First reductions}

We use the abbreviation \((f,F,g)\) for the data of Theorem 5.1, in any one of the three cases considered there.

\begin{statement}{Lemma 6.1}
For Theorem 5.1 to hold for fixed \(f\) and \(F\), and arbitrary \(g\), it is enough that it hold after replacing \(S\) by every strict localization \(\bar S\) of a scheme \(S'\) locally of finite type over \(S\), at a point closed in its fibre over \(S\). More precisely, if \(\bar g:\bar s\to \bar S\) is the inclusion of the closed point and \((\bar f,\bar F)\) is obtained from \((f,F)\) by base change, it is enough to know the theorem for \((\bar f,\bar F,\bar g)\).
\end{statement}

\begin{prooftext}
One may assume \(S=\Spec A\) and \(S'=\Spec A'\). Write \(A'\) as the filtered inductive limit of its finite-type \(A\)-subalgebras \(A_i'\). The limit formalism of Expos\'e VII, 5.11 and 5.14, reduces the question to the case where \(g\) is of finite type. To prove that the base-change homomorphisms are isomorphisms, it is enough to test on geometric fibres at points of \(S'\) which are closed in their fibre over \(S\) (VIII, 3.13). If \(\bar s'\) is such a geometric point, the residue field extension is algebraic, so \(\bar s'\) may also be regarded as a geometric point \(\bar s\) of \(S\). Let \(\bar S\) and \(\bar S'\) be the corresponding strict localizations. There is a commutative square
\[
\xymatrix@C=3cm@R=1.2cm{
\bar s'\ar[r]\ar[d]&\bar s\ar[d]\\
\bar S'\ar[r]&\bar S,
}
\]
whose upper horizontal arrow is an isomorphism. By VIII, 5.2 and 5.3, the map induced on the fibre at \(\bar s'\) by the base-change morphism for \((f,F,g)\) identifies with
\[
H^i(\bar X,\bar F)\longrightarrow H^i(\bar X',\bar F'),
\]
where \(\bar X\) and \(\bar X'\) are obtained by base change from \(\bar S\) and \(\bar S'\). This map fits into the evident square comparing the cohomology of the special fibres. The top arrow is an isomorphism, and the two vertical arrows are isomorphisms by hypothesis. Hence the bottom arrow is an isomorphism. \qedtext
\end{prooftext}

Thus, in proving Theorem 5.1 for a fixed situation \((F,f,-)\), we are reduced to the special case of Corollary 5.5 with \(S\) strictly local. We shall need a general criterion for a morphism \(h:Y\to X\) to induce bijections
\[
H^i(X,F)\longrightarrow H^i(Y,h^*F).
\]
The following elementary statements are recorded in somewhat greater generality than is strictly needed later.

\begin{statement}{Lemma 6.2}
Let \(C\) and \(C'\) be abelian categories, let \(T^\bullet\) and \({T'}^\bullet\) be cohomological functors from \(C\) to \(C'\), and let \(\varphi:T^\bullet\to {T'}^\bullet\) be a morphism of cohomological functors. Assume both functors vanish in negative degrees and that \(T^i\) is effaceable for \(i>0\). For an integer \(n\), the following are equivalent:
\begin{enumerate}[label=\alph*)]
\item \(\varphi^i\) is an isomorphism for \(i\le n\) and a monomorphism for \(i=n+1\).
\item If \(n\ge -1\), \(\varphi^0\) is a monomorphism and \(\varphi^i\) is an epimorphism for \(i\le n\).
\item If \(n\ge0\), \(\varphi^0\) is an isomorphism and \({T'}^i\) is effaceable for \(0<i\le n\).
\end{enumerate}
When these conditions hold, \(\varphi^{n+1}(A):T^{n+1}(A)\to {T'}^{n+1}(A)\) is an epimorphism if and only if \({T'}^{n+1}(A)\) is effaceable. If \(C'=\mathbf{Ab}\) and \(\xi\in {T'}^{n+1}(A)\), then \(\xi\) is in the image of \(T^{n+1}(A)\) if and only if it is effaceable.
\end{statement}

\begin{prooftext}
The implication (a) \(\Rightarrow\) (b) and the implication (a) \(\Rightarrow\) (c) are immediate. Since the positive \(T^i\) are right satellites of \(T^0\), the converse implications are the usual uniqueness and effaceability criterion for universal \(\delta\)-functors. Equivalently, one argues by induction on \(n\), embedding an object \(A\) in an object \(B\) which effaces \(T^{n+1}(A)\), and applying the five lemma to the long exact cohomology sequences associated with \(0\to A\to B\to C\to0\). This also gives the final effaceability criterion. \qedtext
\end{prooftext}

\begin{statement}{Corollary 6.3}
Suppose that \(T^\bullet\) and \({T'}^\bullet\) extend to cohomological functors on \(D^+(C)\), vanishing in negative degrees on complexes whose cohomology vanishes in negative degrees. Let \(K^\bullet\in D^+(C)\) with \(H^i(K^\bullet)=0\) for \(i<0\).
\begin{enumerate}[label=(\roman*)]
\item If condition (a) of Lemma 6.2 holds, then
\[
T^i(K^\bullet)\longrightarrow {T'}^i(K^\bullet)
\]
is an isomorphism for \(i<n\), and a monomorphism for \(i=n+1\).
\item If \(u:L^\bullet\to K^\bullet\) induces an injection on \(H^0\), then surjectivity of
\[
T^{n+1}(L^\bullet)\to {T'}^{n+1}(L^\bullet)
\]
is equivalent to the condition that the image of \({T'}^{n+1}(L^\bullet)\) in \({T'}^{n+1}(K^\bullet)\) lie in the image of \(T^{n+1}(K^\bullet)\). In the case \(C'=\mathbf{Ab}\), this is the corresponding elementwise criterion.
\end{enumerate}
\end{statement}

\begin{prooftext}
For (i), truncate \(K^\bullet\) at the least degree where its cohomology is nonzero and argue by descending induction, using the exact triangle
\[
H^m(K^\bullet)[-m]\longrightarrow K^\bullet\longrightarrow K'{}^\bullet\longrightarrow.
\]
The associated long exact sequences and the five lemma give the assertion. For (ii), use the exact triangle attached to a mapping cylinder of \(u\), and the same diagram chase; the first two vertical arrows are isomorphisms, while the last two are monomorphisms by (i).
\end{prooftext}

\begin{statement}{Lemma 6.4}
Theorem 5.1 holds when \(f\) is finite.
\end{statement}

This is exactly VIII, 5.5 and 5.8.

\begin{statement}{Proposition 6.5}
Let \(h:Y\to X\) be quasi-compact and quasi-separated. In (i), let \(I\) be a set with at least two elements; in (ii) and (iii), let \(\LL\) be a nonempty set of primes; in (iii), let \(n\) be an integer.

\begin{enumerate}[label=(\roman*)]
\item The following are equivalent: for every sheaf of sets \(F\) on \(X\), the map
\[
H^0(X,F)\to H^0(Y,h^*F)
\]
is injective (respectively bijective); equivalently, after every finite morphism \(X'\to X\), the induced map on locally constant maps with value \(I\), or equivalently on clopen subsets, is injective (respectively bijective). If \(X\) is locally noetherian, this can be stated as surjectivity (respectively bijectivity) of \(\Pi_0(Y')\to\Pi_0(X')\) for all finite \(X'\to X\). Under these conditions, pullback on torsors under any sheaf of groups is faithful (respectively fully faithful), and pullback on \(\etale\) coverings is faithful (respectively fully faithful).
\item The following are equivalent: for every sheaf of ind-\(\LL\)-finite groups \(F\), the maps in \(H^0\) and \(H^1\)
\[
H^i(X,F)\to H^i(Y,h^*F),\qquad i=0,1,
\]
are bijective; equivalently this may be tested after every finite \(X'\to X\) on constant finite \(\LL\)-groups. If \(\LL\) is the set of all primes, this is the same as saying that pullback induces an equivalence on categories of \(\etale\) coverings after every finite base change; in the noetherian case, it may be phrased in terms of \(\Pi_0\) and \(\Pi_1\).
\item The following are equivalent: for every \(\LL\)-torsion abelian sheaf \(F\),
\[
H^i(X,F)\to H^i(Y,h^*F)
\]
is an isomorphism for \(i\le n\) and a monomorphism for \(i=n+1\); equivalently, for \(n\ge -1\), it is injective in degree \(0\) and surjective in degrees \(i\le n\); equivalently, this can be tested after finite \(X'\to X\) on the constant sheaves \(\Zz/\ell^\nu\Zz\), \(\ell\in\LL\), \(\nu>0\).
\end{enumerate}
\end{statement}

\begin{statement}{Corollary 6.6}
Assume the hypotheses of Proposition 6.5. Then the following lifting criteria hold.
\begin{enumerate}[label=(\roman*)]
\item Under the injective form of 6.5(i), if \(\xi\in H^0(Y,h^*F)\) becomes induced from \(X\) after a monomorphism \(F\hookrightarrow G\), then \(\xi\) was already induced from \(X\).
\item Under the bijective form of 6.5(i), the same assertion holds for classes \(\xi\in H^1(Y,h^*F)\) of torsors under a sheaf of groups.
\item Under 6.5(iii) for a given \(n\), the same assertion holds for \(\xi\in H^{n+1}(Y,h^*F)\), where \(F\) is an abelian \(\LL\)-torsion sheaf.
\end{enumerate}
\end{statement}

\begin{prooftext}
The abelian assertion is Lemma 6.2 (or Corollary 6.3 in the derived formulation) applied to sheaves of \(\LL\)-torsion. For (i), form the pushout \(H=G\amalg_FG\); since monomorphisms in a topos are effective, \(F\to G\rightrightarrows H\) is exact, and a diagram chase in global sections gives the claim. For (ii), use the sheaf \(\Isom_F(P,P')\): the hypothesis on \(H^0\) implies that pullback on \(F\)-torsors is faithful or fully faithful. If a \(h^*F\)-torsor becomes the pullback of a \(G\)-torsor, it is described by a section of the quotient torsor \(Q/F\); the bijectivity of \(H^0\) lifts this section to \(X\), hence descends the reduction of structure group.

For Proposition 6.5 itself, reduce constructible sheaves by IX, 2.7.2 and IX, 2.14 to finite direct images of constant sheaves along finite maps. In the abelian case use the commutation of cohomology with filtered inductive limits (VII, 3.3) and Lemma 6.2. In the non-abelian case the same reduction uses effaceability of non-abelian \(H^1\) and the exact sequence of torsors. The interpretations by clopen subsets, connected components, and \(\etale\) coverings are the standard Galois-theoretic translations.
\end{prooftext}

\begin{remarktext}{Remark 6.7}
The arguments proving 6.5 and 6.6 are formal. They could be isolated as abstract lemmas in which the sheaves \(p_*G_{X'}\), for finite \(p:X'\to X\), play the role of cogenerators. When \(X\) is noetherian, the tests in 6.5 may often be restricted to integral finite \(X'\), and, for universally Japanese \(X\), even to normal integral \(X'\). The proof of 6.6 and the implications from (a) to the other forms do not use quasi-compactness or quasi-separatedness; those hypotheses enter only through passage to filtered limits of sheaves.
\end{remarktext}

\begin{statement}{Proposition 6.8 (descent lemma)}
With the notation of Proposition 6.5, let \(p:\bar X\to X\) be a morphism, let \(\bar Y=Y\times_X\bar X\), and let \(\bar h:\bar Y\to\bar X\), \(p_Y:\bar Y\to Y\) be the evident maps. Assume either that \(p\) is surjective, or that \(h\) is a closed immersion and \(X=p(\bar X)\cup h(Y)\).
\begin{enumerate}[label=(\roman*)]
\item If \(\bar h\) satisfies the injective (respectively bijective) condition of 6.5(i), and in the bijective case if
\[
h^*p_*\bar F\longrightarrow p_{Y*}\bar h^*\bar F
\]
is injective for every sheaf of sets \(\bar F\) on \(\bar X\), then \(h\) satisfies the corresponding condition 6.5(i).
\item If \(\bar h\) satisfies 6.5(ii) and the same base-change morphism is bijective for every ind-\(\LL\)-group sheaf \(\bar F\), then \(h\) satisfies 6.5(ii).
\item If \(\bar h\) satisfies 6.5(iii), and, for every abelian \(\LL\)-torsion sheaf \(\bar F\) on \(\bar X\),
\[
h^*R^ip_*\bar F\longrightarrow R^ip_{Y*}\bar h^*\bar F
\]
is bijective for \(i\le n-1\) and injective for \(i=n\), then \(h\) satisfies 6.5(iii).
\end{enumerate}
\end{statement}

\begin{prooftext}
If necessary replace \(\bar X\) by the disjoint union \(\bar X\amalg Y\), so that \(p\) is surjective. Then every sheaf \(F\) on \(X\) injects into \(p_*p^*F\). In case (i), reduce to sheaves of the form \(p_*\bar F\) and compare the square of global sections on \(X,Y,\bar X,\bar Y\); injectivity or surjectivity follows from the hypothesis on \(\bar h\) and the base-change map. In case (ii), use the non-abelian exact sequence for \(H^1\) attached to \(p\) and \(p_Y\); the middle vertical map is bijective by the hypothesis on \(\bar h\), while the terms measuring obstruction are controlled by (i) and by the base-change hypothesis. In case (iii), replace the injection \(F\to p_*p^*F\) by the morphism \(F\to Rp_*p^*F\) in \(D^+\), apply Corollary 6.3, and compare hypercohomology through the base-change morphism to \(Rp_{Y*}\bar h^*\bar F\). The composite is the map induced by \(\bar h\), hence an isomorphism in the prescribed range.
\end{prooftext}

\begin{statement}{Corollary 6.9}
In the situation of 6.8, assume in addition that \(h\) itself satisfies the relevant condition of 6.5 and that the base-change morphisms for \(p\) are isomorphisms in the same range. Let \(F\) be a sheaf on \(X\) of the corresponding type, and let \(\xi\) be a class on \(Y\). Then \(\xi\) comes from \(X\) if and only if its pullback to \(\bar Y\) comes from \(\bar X\).
\end{statement}

This is exactly the sharpened form of the proof of 6.8: the hypercohomology comparison shows that the only obstruction is detected after the auxiliary cover \(p\).

\begin{statement}{Corollary 6.10}
Assume the injective form of 6.5(i). To obtain the bijective form of 6.5(i), or 6.5(ii), or 6.5(iii), it is necessary and sufficient that the following local lifting condition hold. For every sheaf \(F\) of the relevant type, every class \(\xi\) in the relevant cohomology of \(Y\), and every nonempty closed subset \(X'\subset X\), there exists a morphism \(p:\bar X\to X\) such that:
\begin{enumerate}[label=\arabic*\textdegree)]
\item \(p(\bar X)\subset X'\) and contains a nonempty open subset of \(X'\);
\item the relevant base-change morphisms for \(p\) are isomorphisms in the necessary degrees;
\item the pullback of \(\xi\) to \(\bar Y\) lies in the image of the cohomology of \(\bar X\).
\end{enumerate}
Moreover, if the theorem is already known through degree \(n\), the same criterion characterizes whether a class in degree \(n+1\) descends.
\end{statement}

\begin{prooftext}
One direction is immediate from 6.9. Conversely, suppose a class \(\xi\) does not descend. Let \(\Phi\) be the set of closed subsets \(X'\subset X\) on which the restriction of \(\xi\) still does not descend, ordered by reverse inclusion. The limit theorem VII, 5.8, shows that \(\Phi\) is inductive; choose a minimal element. Applying the stated local lifting condition to this minimal closed set, and enlarging \(p\) by adjoining the complementary closed part if needed, we may assume \(p\) is surjective over it. Corollary 6.9 then forces descent on the minimal closed set, a contradiction.
\end{prooftext}

\begin{statement}{Proposition 6.11 (transitivity lemma)}
With the notation of 6.8, but without assuming \(p\) surjective, suppose that \(h\) satisfies the relevant condition of 6.5. Suppose also that the corresponding base-change morphisms for \(p\) are monomorphisms or isomorphisms in the same range: for sheaves of sets in case (i), for ind-\(\LL\)-groups in case (ii), and for abelian \(\LL\)-torsion sheaves in case (iii). Then \(\bar h:\bar Y\to\bar X\) satisfies the same condition.
\end{statement}

\begin{prooftext}
For sheaves of sets, the diagram used in the proof of 6.8 is read backwards. For ind-\(\LL\)-groups, efface the two obstruction terms in the exact sequence for \(H^1\) attached to \(p\), using Proposition 3.3; the base-change hypotheses identify these terms with the corresponding terms over \(X\), where they are effaceable. In the abelian case, compare
\[
R\Gamma(\bar X,\bar F)\longrightarrow R\Gamma(\bar Y,\bar h^*\bar F)
\]
with the composite
\[
R\Gamma(X,Rp_*\bar F)\longrightarrow R\Gamma(Y,h^*Rp_*\bar F)
\longrightarrow R\Gamma(Y,Rp_{Y*}\bar h^*\bar F).
\]
The first arrow is controlled by the hypothesis on \(h\) and Corollary 6.3, and the second by the base-change hypothesis. \qedtext
\end{prooftext}

\begin{remarktext}{Remarks 6.12--6.13}
The proof in cases (i) and (iii) gives a more precise statement for a fixed sheaf \(\bar F\): it is enough to know base change for \(R^ip_*\bar F\) and the corresponding cohomological comparison for each \(R^jp_*\bar F\). An analogous fixed-sheaf statement for case (ii) would involve non-abelian two-dimensional cohomology, which is why the proof above uses effaceability instead.

The principal example of the conditions in 6.5 is the situation of Corollary 5.5, proved in the next expos\'e using the reductions of the present one. Similar questions arise for henselian pairs \((X,Y)\), for formal completions defined by an ideal, and for putative ``\'etale localizations along a closed subset''. One may conjecture that henselian pairs satisfy the non-abelian and abelian forms of 6.5, and ask when a weaker condition on sections of \(\etale\) covers already implies henselianity. These remarks anticipate later developments and will not be used in the proof.
\end{remarktext}

\section*{7. A variant of Chow's lemma}

We need the following form of Chow's lemma, stripped of noetherian and irreducibility hypotheses.

\begin{statement}{Lemma 7.1}
Let \(S\) be quasi-compact and quasi-separated, and let \(f:X\to S\) be separated of finite type, with \(X\ne\varnothing\). There is a commutative diagram
\[
\xymatrix@=1.25cm{
X\ar[rd]_f&&\bar X\ar[ll]_{\pi}\ar[ld]^{\bar f}\\
&S&
}
\]
where \(\pi\) is projective, \(\bar f\) is quasi-projective, and there is a nonempty open \(U\subset X\) such that \(\pi\) induces an isomorphism \(\pi^{-1}(U)\simeq U\).
\end{statement}

\begin{prooftext}
First note a permanence property: if the lemma holds for a closed subscheme \(X'\subset X\) and for a dense open \(V\subset X'\) which is also open in \(X\), then it holds for \(X\), by composing the projective morphism \(\bar X'\to X'\) with the closed immersion \(X'\to X\). Cover \(X\) by finitely many nonempty affine opens \(U_1,
\ldots,U_n\). Since \(S\) is quasi-compact and quasi-separated, each \(U_i\to S\) is quasi-affine. If \(U=\bigcap U_i\) is dense, the usual proof of Chow's lemma (EGA II, 5.6.1) applies. If the intersection is nonempty but not dense, replace \(X\) by the closure of that intersection. If the intersection is empty, choose a maximal nonempty partial intersection and replace \(X\) by its closure, reducing the number of open sets. Induction on the number of opens gives the result.
\end{prooftext}

\section*{8. Final reductions}

\begin{statement}{Lemma 8.1}
To prove any of the assertions 5.1(i), (ii), (iii), it is enough to prove it when \(f\) is projective and \(S\) is noetherian.
\end{statement}

\begin{prooftext}
Assume first that the theorem is known for projective \(f\). By Lemma 6.1, it is enough to handle the case of Corollary 5.5 with \(S\) strictly local. Apply Corollary 6.10 with \(Y=X_0\), the closed fibre. Given the closed set and class occurring there, use Lemma 7.1 to find \(p:\bar X\to X\) whose image contains a nonempty open part and with \(\bar X\) projective over \(S\). The projective case supplies the required base-change isomorphisms and the cohomological lifting condition, so Corollary 6.10 gives the theorem.

It remains to reduce the projective case to noetherian base. Again by 6.1 we may take \(S\) strictly local and prove Corollary 5.5. Embed \(X\) in \(\mathbb P^r_S\) and replace \(X\) by \(\mathbb P^r_S\) using the compatibility already established in 4.4. If \(S=\Spec A\), write \(A\) as a filtered limit of noetherian henselian local rings \(A_i\) mapping locally to \(A\), and use VII, 5.7 to commute cohomology with this limit. The closed fibres are obtained by the corresponding filtered limit of residue fields. Thus the theorem over \(A\) follows from the theorem over the noetherian \(A_i\).
\end{prooftext}

\begin{statement}{Corollary 8.2}
Theorem 5.1(i) is true.
\end{statement}

Indeed, combine Lemma 8.1, the criterion 6.5(i) in the connected-component form, and Lemma 5.8.

\begin{statement}{Lemma 8.3}
To prove any of 5.1(i), (ii), (iii), it is enough to prove it for \(f\) projective of relative dimension \(\le1\), with \(S\) noetherian.
\end{statement}

\begin{prooftext}
By Lemma 8.1 we reduce to \(X=\mathbb P^r_S\). Induct on \(r\). For \(r=1\) this is the stated relative-dimension case. For \(r\ge2\), blow up the standard \(\mathbb P^{r-2}_S\subset\mathbb P^r_S\). The blow-up \(\bar X\) maps both to \(X\) and to \(X_1=\mathbb P^1_S\); over \(X_1\) it is locally a projective \((r-1)\)-space. The fibres of \(\bar X\to X\) and \(X_1\to S\) have dimension at most one. Applying the dimension-one hypothesis to the former and to \(X_1\to S\), and the induction hypothesis to \(\bar X\to X_1\), the descent and transitivity lemmas 6.8 and 6.11 give the result for \(X\to S\).
\end{prooftext}

\begin{statement}{Lemma 8.4}
To prove Theorem 5.1(ii), it is enough to prove the specialization theorem for the fundamental group, Theorem 5.9 bis, when \(f\) is projective of relative dimension \(\le1\) and \(S\) is noetherian strictly local.
\end{statement}

This is immediate from Lemma 8.3, Lemma 6.1, and criterion 6.5(ii).

\begin{statement}{Lemma 8.5}
To prove Theorem 5.1(iii), it is enough to prove Theorem 5.1(ii) and the following Picard-surjectivity statement.
\end{statement}

\begin{statement}{Proposition 8.6}
Let \(S\) be henselian noetherian, let \(f:X\to S\) be projective of relative dimension \(\le1\), and let \(X_0\) be the closed fibre. Then
\[
\Pic X\longrightarrow \Pic X_0
\]
is surjective.
\end{statement}

\begin{prooftext}
By 6.5(iii) and Corollary 8.2, it remains to prove surjectivity of
\[
H^q(X,\Zz/\ell^\nu\Zz)\longrightarrow H^q(X_0,\Zz/\ell^\nu\Zz)
\]
for \(q\ge1\). Degree \(1\) follows from 5.1(ii). Since \(X_0\) has dimension at most one over a separably closed field, the groups vanish for \(q>2\) if \(\ell\ne\operatorname{char}k\), and for \(q>1\) if \(\ell=\operatorname{char}k\). It remains to handle \(q=2\), \(\ell\ne\operatorname{char}k\). Put \(n=\ell^\nu\). Since \(S\) is strictly local, \(\mu_n\simeq\Zz/n\Zz\) on \(S\), hence on \(X\). The Kummer sequences for \(X\) and \(X_0\) give a diagram
\[
\xymatrix@R=1.2cm{
\Pic X\ar[r]\ar[d]&H^2(X,\mu_n)\ar[r]\ar[d]&H^2(X,\Gm)\ar[d]\\
\Pic X_0\ar[r]&H^2(X_0,\mu_n)\ar[r]&H^2(X_0,\Gm).
}
\]
Since \(H^2(X_0,\Gm)=0\) by IX, 4.6, surjectivity of \(\Pic X\to\Pic X_0\) implies surjectivity in degree \(2\), proving the lemma.
\end{prooftext}

\begin{thebibliography}{99}
\bibitem{XII-Giraud} J. Giraud, \emph{Non-abelian cohomology}, thesis and Comptes Rendus notes.
\bibitem{XII-Godement} R. Godement, \emph{Th\'eorie des faisceaux}, Paris, 1958.
\bibitem{XII-Verdier} J.-L. Verdier, \emph{Cat\'egories d\'eriv\'ees}, IHES, 1964.
\end{thebibliography}

% ============================================================
\clearpage
\begin{center}
{\Large \textbf{SGA 4, Expos\'e XIII}}\\[0.35em]
{\large \textbf{The Proper Base Change Theorem: End of the Proof}}\\[0.25em]
{\normalsize M. Artin}
\end{center}

\section*{1. The projective flat case}

We recall Theorem 5.9 bis of Expos\'e XII in the case needed here.

\begin{statement}{Proposition 1.1}
Let \(S\) be the spectrum of a noetherian henselian ring, let \(f:X\to S\) be projective and flat, and let \(X_0\) be the closed fibre. Then the restriction functor
\[
\Et(X)\longrightarrow \Et(X_0)
\]
is an equivalence of categories.
\end{statement}

Full faithfulness follows already from XII, 6.5(i). It remains to show that every \(\etale\) covering of \(X_0\) comes from one of \(X\).

\begin{statement}{Lemma 1.2}
Let \(S\) be locally noetherian, let \(f:X\to S\) be projective and flat, and let \(Y\subset X\) be a subscheme. The subfunctor of the final functor on \((\Sch)/S\) defined by the condition \(Y'=X'\) after base change to \(S'\) is representable by a subscheme \(Z\subset S\). If \(Y\) is open (respectively closed), then so is \(Z\).
\end{statement}

\begin{prooftext}
If \(Y\) is closed in an open \(U\subset X\), the condition \(U'=X'\) is represented by the complement of the proper image of \(X-U\). We reduce to \(Y\) closed. Let \(\Oo_X\twoheadrightarrow\Oo_Y\). The condition \(Y=X\) is equivalent to the existence of a splitting \(\Oo_Y\to\Oo_X\), or equivalently to lifting the identity section of \(\Hhom(\Oo_X,\Oo_X)\) to a section of \(\Hhom(\Oo_Y,\Oo_X)\). By EGA III, 7.7.8, these section functors are represented by vector bundles over \(S\), and the inverse image of the identity section is a closed subscheme of \(S\).
\end{prooftext}

\begin{statement}{Lemma 1.3}
Let \(S\) be locally noetherian, \(f:X\to S\) projective and flat, and \(E\) locally free on \(X\). The functor \(QE\), assigning to \(S'/S\) the quotients \(A'\) of \(E'\) endowed with the structure of an \(\Oo_{X'}\)-algebra \(\etale\) over \(X'\), is represented by a scheme locally of finite type over \(S\).
\end{statement}

\begin{prooftext}
The relative \(\Quot(E)\) functor is represented by a scheme locally of finite type. The condition that the universal quotient be flat over \(X'\) is open, so we reduce to the following representability statement.
\end{prooftext}

\begin{statement}{Lemma 1.4}
Let \(S\) be locally noetherian, \(f:X\to S\) projective and flat, and \(M\) locally free on \(X\). The functor assigning to \(S'/S\) the algebra structures on \(M'\) which are \(\etale\) over \(X'\) is represented by a scheme of finite type over \(S\).
\end{statement}

\begin{prooftext}
Bilinear multiplication laws on \(M\) are represented by the vector bundle associated with \(\Hhom(M\otimes M,M)\). The associativity and commutativity conditions are equalities of sections of suitable vector bundles; Lemma 1.2 represents the locus where they hold. After imposing these, the existence of a unit is represented by a vector bundle of sections of \(M\), again cut out by an equality of sections. Finally, for a commutative unital algebra structure, the condition of being \(\etale\) is open. This proves Lemma 1.4 and hence Lemma 1.3.
\end{prooftext}

\subsection*{1.5}
Let \(Z/S\) represent the functor in Lemma 1.3. We examine smoothness of \(Z\) over \(S\). Translating SGA 1, III, 3.1, into the language of \(QE\), one obtains the following infinitesimal lifting criterion. Let \(B\) be an artinian local ring with maximal ideal \(\mathfrak m\), let \(J\subset B\) satisfy \(J\mathfrak m=0\), let \(S'=
\Spec B\) and \(S''=\Spec B/J\). Given a quotient \(A''\) of \(E''\) with an \(\etale\) algebra structure over \(X''\), it should lift to a quotient \(A'\) of \(E'\) with \(\etale\) algebra structure over \(X'\).

Since \(X'\) and \(X''\) have the same underlying space, an \(\etale\) algebra \(A'\) lifting \(A''\) is unique up to unique isomorphism. The only obstruction is therefore to lifting the surjection \(E''\twoheadrightarrow A''\). The exact sequence
\[
0\to J\to B\to B/J\to0
\]
gives
\[
0\to J\otimes\Hhom(E',A')\to \Hhom(E',A')\to \Hhom(E'',A'')\to0,
\]
so the obstruction lies in \(H^1(X',J\otimes\Hhom(E',A'))\). Since \(J\mathfrak m=0\), this is \(J\otimes H^1(X_0,\Hhom(E_0,A_0))\), after compatibility of cohomology with base-field extension.

\begin{statement}{Corollary 1.6}
Let \(Z/S\) represent the functor of Lemma 1.3, let \(X_Z=X\times_SZ\), and let \(A_Z\) be the universal \(\etale\) quotient algebra. Then \(Z\) is smooth over \(S\) at every point \(z_0\) for which
\[
H^1(X_0,\Hhom(E_0,A_0))=0,
\]
where the subscript \(0\) denotes restriction to the fibre of \(X_Z/Z\) over \(z_0\).
\end{statement}

We now finish the proof of Proposition 1.1. Let \(Y_0\to X_0\) be an \(\etale\) covering, corresponding to an \(\etale\) algebra \(A_0\). For \(n\) large, \(A_0(n)\) is generated by global sections, hence there is a surjection
\[
\Oo_{X_0}^{N}(-n)=E_0\twoheadrightarrow A_0.
\]
Choose \(n\) even larger so that
\[
H^1(X_0,\Hhom(E_0,A_0))=0.
\]
Let \(E=\Oo_X^{N}(-n)\), and let \(Z/S\) be the representing scheme. The quotient \(A_0\) determines a rational point \(z_0\) of the closed fibre of \(Z\), and by Corollary 1.6, \(Z\) is smooth at \(z_0\). Since \(S\) is henselian, a smooth \(S\)-scheme with a point in the closed fibre has a section through that point. The corresponding quotient of \(E\) is an \(\etale\) algebra on \(X\) lifting \(A_0\), proving essential surjectivity.

\section*{2. The case of relative dimension \(\le1\)}

\begin{statement}{Proposition 2.1}
With the notation of XII, 5.9 bis, if \(S\) is noetherian strictly local and \(f:X\to S\) is projective of relative dimension \(\le1\), then
\[
\Et(X)\longrightarrow \Et(X_0)
\]
is an equivalence of categories. Together with XII, 8.5, this proves XII, 5.1(ii).
\end{statement}

\begin{statement}{Lemma 2.2}
Let \(S=\Spec A\) be local noetherian, and let \(X/S\) be projective of relative dimension \(\le n\). There exists a projective flat \(S\)-scheme \(Z\) of relative dimension \(\le n\), and a closed immersion \(X\hookrightarrow Z\) over \(S\).
\end{statement}

\begin{prooftext}
Embed \(X\) in \(\mathbb P^N_S\). If \(J\) is the homogeneous ideal defining \(X\), then over the closed fibre, of dimension \(\le n\), choose \(N-n\) homogeneous equations in the image of \(J\) cutting out a subscheme of dimension \(n\). Lift them to elements of \(J\), and let \(Z\) be the closed subscheme they define in \(\mathbb P^N_S\). By EGA IV, 11.3, \(Z\) is flat over \(S\), and by EGA IV, 13.1.5, its relative dimension is \(n\). \qedtext
\end{prooftext}

\begin{prooftext}[Proof of Proposition 2.1]
Assume \(X\) and \(X_0\) connected. Embed \(X\) as a closed subscheme of a projective flat \(Z/S\) of relative dimension at most one. We have
\[
\xymatrix@C=2.5cm@R=1.2cm{
\Et(Z)\ar[r]\ar[d]&\Et(X)\ar[d]\\
\Et(Z_0)\ar[r]&\Et(X_0).
}
\]
The right vertical functor is fully faithful by XII, 5.8. It remains to show essential surjectivity. The left vertical functor is essentially surjective by Proposition 1.1. Thus it is enough to show that every \(\etale\) covering of \(X_0\) is induced from one of \(Z_0\). We are reduced to the following proposition over a separably closed field.
\end{prooftext}

\begin{statement}{Proposition 2.3}
Let \(k\) be separably closed, let \(Z\) be locally of finite type over \(k\) and of dimension \(\le1\), and let \(X\hookrightarrow Z\) be a connected closed subscheme. Every \(\etale\) covering of \(X\) is induced by an \(\etale\) covering of \(Z\).
\end{statement}

\begin{prooftext}
The zero-dimensional case is trivial. Otherwise write \(Z=X\cup Y\), where \(Y\) is the closure of \(Z-X\), and put \(V=X\cap Y\), a discrete locally artinian scheme. The finite surjective morphism \(X\amalg Y\to Z\) is an effective descent morphism for \(\etale\) coverings. Moreover
\[
(X\amalg Y)\times_Z(X\amalg Y)=X\amalg Y\amalg V^{(1)}\amalg V^{(2)},
\]
where the two copies of \(V\) are exchanged between the \(X\)- and \(Y\)-components by the projections. Given an \(\etale\) covering \(X'\to X\) of degree \(n\), choose any \(\etale\) covering \(Y''\to Y\) of degree \(n\), for instance the trivial one. A descent datum is just an isomorphism between the two induced covers of \(V\). Since \(k\) is separably closed, every \(\etale\) covering of \(V\) is split, so such an isomorphism exists.
\end{prooftext}

\begin{remarktext}{Remark 2.4}
One can avoid the use of the full descent theorem by proving directly that \(\Et(Z)\) is equivalent to the category of triples \((X',Y',\theta)\), where \(X'\) and \(Y'\) are \(\etale\) coverings of \(X\) and \(Y\), and \(\theta\) identifies their pullbacks to \(V\). This follows from the analogous gluing statement for quasi-coherent modules on a union of two closed subschemes.
\end{remarktext}

\section*{3. An auxiliary result on the Picard group}

\begin{statement}{Proposition 3.1}
Let \(X\) be noetherian, and let \(X_0\subset X\) be closed such that: (a) \(\dim X_0\le1\); (b) for every nonempty connected closed subset \(Y\subset X\), the intersection \(Y_0=Y\cap X_0\) is nonempty and connected. Let \(Z_0\) be the union, over \(x\in\Ass\Oo_X\), of the sets \(\overline{\{x\}}\cap X_0\) when this intersection is reduced to one point.
\begin{enumerate}[label=(\roman*)]
\item Every Cartier divisor \(D_0\) on \(X_0\) whose support avoids \(Z_0\) is induced by a Cartier divisor \(D\) on \(X\). If \(D_0\ge0\), then \(D\) may be chosen \(\ge0\).
\item If there is an affine open \(U_0\subset X_0\) containing the finite set \(\Ass\Oo_{X_0}\cup Z_0\) (for instance if \(X_0\) has an ample invertible sheaf), then \(\Pic X\to\Pic X_0\) is surjective.
\end{enumerate}
\end{statement}

\begin{prooftext}
For (i), decompose any Cartier divisor on the one-dimensional scheme \(X_0\) into the difference of effective Cartier divisors. At a point \(x\) of the support of an effective divisor, choose a local equation \(f_x\in\Oo_{X_0,x}\) which is a non-zero-divisor and lift it to a section \(g_x\) of \(\Oo_X\) on a neighbourhood \(U_x\). By avoiding the associated points of \(\Oo_X\), using the definition of \(Z_0\) and the connectedness hypothesis, shrink \(U_x\) so that \(g_x\) is a non-zero-divisor. The closure of \(V(g_x)\) has a connected component meeting \(X_0\) only at \(x\); that component defines a Cartier divisor on \(X\) inducing the desired local contribution. Summing over the finite support gives \(D\).

For (ii), every invertible sheaf on \(X_0\) is represented by a Cartier divisor whose support avoids \(Z_0\): on the affine open containing \(\Ass\Oo_{X_0}\cup Z_0\), the invertible sheaf is trivial after passing to the corresponding semilocal ring. Apply (i).
\end{prooftext}

\begin{statement}{Corollary 3.2}
Let \(S\) be the spectrum of a noetherian henselian ring, and let \(f:X\to S\) be proper of relative dimension \(\le1\). If \(X'_0\subset X\) is a closed subscheme with the same underlying space as the closed fibre \(X_0\), then
\[
\Pic X\longrightarrow \Pic X'_0
\]
is surjective.
\end{statement}

Apply Proposition 3.1 to \((X,X'_0)\). The dimension condition is immediate, and connectedness of intersections follows from XII, 5.7. The affine-open hypothesis follows because a proper curve over a field is projective, hence has ample invertible sheaves. This completes the proof of XII, 5.1(iii).

\begin{remarktext}{Remark}
If one does not wish to use projectivity of proper curves over a field, Corollary 3.2 may be used only in the projective case, which is enough for the application in XII. Conversely, accepting projectivity of proper algebraic curves, Corollary 3.2 also implies that \(X\) is projective over \(S\) under its hypotheses: lift an ample class from \(X_0\) to \(X\) and apply EGA III, 4.7.1.
\end{remarktext}

\begin{thebibliography}{99}
\bibitem{TDTE} A. Grothendieck, \emph{Technique de descente et th\'eor\`emes d'existence en g\'eom\'etrie alg\'ebrique}, in \emph{Fondements de la G\'eom\'etrie Alg\'ebrique}.
\end{thebibliography}

% ============================================================
\clearpage
\begin{center}
{\Large \textbf{SGA 4, Expos\'e XIV}}\\[0.35em]
{\large \textbf{A Finiteness Theorem for a Proper Morphism; Cohomological Dimension of Affine Algebraic Schemes}}\\[0.25em]
{\normalsize M. Artin}
\end{center}

This expos\'e contains two important theorems whose proofs use the proper base change theorem.

\section*{1. Finiteness theorem for a proper morphism}

\begin{statement}{Theorem 1.1}
Let \(f:X\to S\) be proper and of finite presentation. Let \(A\) be a left noetherian ring which is a \(\Zz/n\Zz\)-algebra for some \(n\). Let \(F\) be a constructible sheaf of sets (respectively groups, respectively \(A\)-modules) on \(X\). Then \(f_*F\) (respectively \(R^1f_*F\), respectively \(R^qf_*F\) for every \(q\ge0\)) is constructible.
\end{statement}

\begin{statement}{Corollary 1.2}
If \(X\) is proper over the spectrum of a separably closed field \(k\), and \(F\) is a constructible torsion abelian sheaf on \(X\), then the groups \(H^q(X,F)\) are finite for all \(q\ge0\).
\end{statement}

\begin{remarktext}{Remark 1.3}
Properness is essential for abelian sheaves: take \(F=\Zz/p\Zz\) on \(\mathbb A^1_k\), with \(p=\operatorname{char}k\), and either compactify into \(\mathbb P^1_k\) or take the structural morphism. One expects broader finiteness results for coefficients prime to the residual characteristics over excellent bases; in characteristic zero this follows from resolution of singularities.
\end{remarktext}

\begin{prooftext}[Proof of Theorem 1.1]
The question is local on \(S\). Write affine \(S\) as an inverse limit of schemes \(S_0\) of finite type over \(\Spec\Zz\). The data descend to some \(S_0\), and proper base change for \(S\to S_0\) reduces constructibility to the case \(S\) of finite type over \(\Spec\Zz\).
\end{prooftext}

\begin{statement}{Lemma 1.4}
It is enough to prove Theorem 1.1 when \(S\) is of finite type over \(\Spec\Zz\).
\end{statement}

\begin{statement}{Lemma 1.5}
The assertion of Theorem 1.1 for sheaves of sets is true: if \(F\) is constructible, then \(f_*F\) is constructible.
\end{statement}

\begin{prooftext}
Assume \(S\) noetherian. Embed \(F\) into a finite product \(G=\prod_i\pi_{i*}C_i\), with \(\pi_i:X_i\to X\) finite and \(C_i\) constant constructible (IX, 2.14). It is enough to treat \(F=C_X\) constant finite. By proper base change, the geometric fibre of \(f_*F\) at \(\bar s\) is \(C^{\pi_0(X_{\bar s})}\). This is finite, and the function \(s\mapsto \#\pi_0(X_{\bar s})\) is constructible by EGA IV, 9.7.9. Hence \(f_*F\) is constructible.
\end{prooftext}

\begin{statement}{Lemma 1.6}
Let \(S\) be noetherian, \(f:X\to S\) proper, and \(F\) constructible, either as a sheaf of groups or of \(A\)-modules. If \(S=\bigcup S_i\) is a finite union of subschemes and the restrictions \(R^qf_{i*}F_i\) are constructible on each \(S_i\), then \(R^qf_*F\) is constructible on \(S\).
\end{statement}

This follows from proper base change, which identifies the restriction of \(R^qf_*F\) to \(S_i\) with \(R^qf_{i*}F_i\), and from the gluing criterion for constructible sheaves (IX, 2.8).

\begin{statement}{Lemma 1.7}
Let \(S\) be noetherian and suppose
\[
\xymatrix@=1.25cm{
X\ar[rd]_f&&\bar X\ar[ll]_{\pi}\ar[ld]^{\bar f}\\
&S&
}
\]
is a commutative diagram of proper morphisms. Let \(j:U\hookrightarrow X\) be an open such that \(\bar U=U\times_X\bar X\to U\) is an isomorphism, and let \(Y=X-U\) with its reduced closed structure. If Theorem 1.1 holds for \(\bar f\), for \(Y\to S\), and for \(\pi\), then it holds for \(f\).
\end{statement}

\begin{prooftext}
For modules, use
\[
0\to j_!j^*F\to F\to i_*i^*F\to0,
\]
where \(i:Y\hookrightarrow X\). The term supported on \(Y\) is handled by the induction hypothesis on \(Y\to S\). For \(j_!j^*F\), one has
\[
j_!j^*F\simeq \pi_*\bar j_!\bar j^*\pi^*F,
\]
verified fibrewise by proper base change. Higher direct images under \(\pi\) of \(\bar j_!\bar j^*\pi^*F\) vanish because \(\pi\) is an isomorphism over \(U\) and the fibres over \(Y\) carry no support. Leray then identifies the required direct images for \(f\) with those for \(\bar f\), which are constructible. The set and group cases follow from the same d\'evissage, using the corresponding exactness statements for \(H^0\) and \(H^1\).
\end{prooftext}

\begin{statement}{Lemma 1.8}
Assume \(S\) is of finite type over \(\Spec\Zz\). It is enough to prove Theorem 1.1 when \(f\) is projective.
\end{statement}

\begin{prooftext}
Apply the Chow lemma of XII, 7.1 to a nonempty closed part of \(X\), and then Lemma 1.7. Noetherian induction on closed subsets of \(X\) gives the result.
\end{prooftext}

\begin{statement}{Lemma 1.9}
It is enough to prove Theorem 1.1 when \(S\) is of finite type over \(\Spec\Zz\) and \(X=\mathbb P^r_S\).
\end{statement}

For a projective \(X\), embed it in \(\mathbb P^r_S\); extend a constructible sheaf by zero, and use the compatibility of direct image with closed immersions.

\begin{statement}{Lemma 1.10}
Theorem 1.1 is true.
\end{statement}

\begin{prooftext}
We now argue by induction on \(r\) for \(X=\mathbb P^r_S\). The case \(r=0\) is finite. For \(r\ge1\), stratify \(S\) so that the constructible sheaf is locally constant on a suitable stratification of \(\mathbb P^r_S\), and apply the hyperplane/blow-up reduction already used in XII, 8.3: blowing up a standard \(\mathbb P^{r-2}\subset\mathbb P^r\) reduces to morphisms of relative dimension \(1\) and to a projective bundle of smaller relative projective dimension. The proper base change theorem controls all restrictions to strata, while Lemma 1.6 glues the constructibility. This completes the proof.
\end{prooftext}

\begin{remarktext}{Remark 1.11}
The proof above is deliberately parallel to the reductions in Expos\'e XII. The noetherian hypotheses enter only through constructibility and noetherian induction; the conclusion is then transported back to arbitrary bases by the limit argument at the start.
\end{remarktext}

\section*{2. A variant of dimension}

For a sheaf \(F\) on a scheme of finite type over a field, let \(d(F)\) denote the smallest integer \(d\) such that the support of \(F\) is contained in a closed subset of dimension \(\le d\); equivalently, after replacing \(F\) by a constructible subsheaf when needed, \(d(F)\) measures the dimension of its support.

\begin{statement}{Definition 2.2}
For a scheme \(X\) locally of finite type over a field, put
\[
\delta(X)=\sup\{\operatorname{trdeg}_{k} k(x)\mid x\in X\},
\]
where \(k\) is the ground field. For a torsion sheaf \(F\), set \(\delta(F)=\delta(\Supp F)\).
\end{statement}

\begin{statement}{Proposition 2.3}
Let \(f:X\to Y\) be a morphism of schemes locally of finite type over a field. For a torsion sheaf \(F\) on \(X\), one has
\[
\delta(R^qf_*F)\le \delta(F)-q
\]
in the range where the cohomological dimension estimates below apply.
\end{statement}

\begin{prooftext}
By constructible approximation and the finiteness theorem, reduce to constructible \(F\). The statement is local on \(Y\) and is checked on strict localizations. The higher direct image at a geometric point is the cohomology of the corresponding fibre, whose dimension is controlled by \(\delta(F)\). The assertion is then the standard dimension drop in cohomology with torsion coefficients.
\end{prooftext}

\begin{statement}{Proposition 2.4}
If \(X\) is affine and algebraic over a separably closed field, and if \(F\) is a torsion sheaf with \(\delta(F)\le d\), then \(H^q(X,F)=0\) for \(q>d\), once the theorem is known in dimensions \(<d\). This proposition is the induction form used in Section 4.
\end{statement}

\section*{3. Cohomological dimension of affine algebraic schemes}

\begin{statement}{Theorem 3.1}
Let \(X\) be an affine algebraic scheme over a separably closed field, and let \(F\) be a torsion abelian sheaf on \(X\). If \(d(F)\le d\), then
\[
H^q(X,F)=0\qquad(q>d).
\]
In particular \(\operatorname{cd}(X)\le\dim X\).
\end{statement}

This is the affine Lefschetz theorem in the form needed here.

\begin{statement}{Corollary 3.2}
If \(X\) is affine of dimension \(d\) over a separably closed field, then every torsion sheaf on \(X\) has vanishing cohomology in degrees \(>d\).
\end{statement}

\begin{statement}{Corollary 3.3}
For any field \(k\), and \(X\) affine algebraic over \(k\), the cohomology of \(X\) is controlled by the cohomological dimension of \(k\) and the dimension of \(X\), via the Hochschild-Serre spectral sequence after base change to a separable closure.
\end{statement}

\begin{statement}{Corollary 3.4}
Let \(f:X\to Y\) be an affine morphism of schemes algebraic over a separably closed field. If \(F\) is torsion and \(d(F)\le d\), then
\[
R^qf_*F=0\qquad(q>d),
\]
and more precisely the support dimension of \(R^qf_*F\) is at most \(d-q\).
\end{statement}

\begin{statement}{Corollary 3.5}
If \(X\) is strictly local and is the strict localization of an algebraic scheme of dimension \(\le d\), and if \(U\subset X\) is the complement of a principal divisor, then \(\operatorname{cd}(U)\le d\).
\end{statement}

\begin{remarktext}{Remark 3.6}
The formulation is stable under passage to strict localizations and is designed for induction in Section 4.
\end{remarktext}

\section*{4. Proof of Theorem 3.1}

\begin{statement}{Lemma 4.1}
The theorem is clear for \(d=0\), and the case \(d=1\) follows from the already established cohomological dimension of affine curves.
\end{statement}

Indeed, if \(d=0\), the support is finite and VIII, 5.5 applies. If \(d=1\), reduce to a constructible sheaf on a scheme of dimension at most one and use IX, 5.7.

\begin{statement}{Lemma 4.2}
To prove Corollary 3.4 in dimension \(d\), it is enough to treat the structural morphism \(\mathbb A^1_{Y'}\to Y'\).
\end{statement}

\begin{prooftext}
Embed the affine morphism in some affine space \(\mathbb A^N_{Y'}\). Induct on \(N\). For \(N>1\), factor the projection as
\[
\mathbb A^N_{Y'}\to \mathbb A^{N-1}_{Y'}\to Y'.
\]
The Leray spectral sequence, together with the induction hypothesis on \(N\) and on support dimension, reduces the assertion to the case of relative affine line.
\end{prooftext}

\begin{statement}{Assertion 4.3(d)}
Let \(X'\) be strictly local of dimension \(\le d\), the strict localization of an algebraic scheme over a field. Let \(f\in\Gamma(X',\Oo_{X'})\), and put \(U=X'-V(f)\). Then \(\operatorname{cd}(U)\le d\).
\end{statement}

\begin{statement}{Lemma 4.4}
Assertion 4.3(d) implies Corollary 3.4(d).
\end{statement}

\begin{prooftext}
Compactify \(\mathbb A^1_{Y'}\) into \(\mathbb P^1_{Y'}\) by adding the section at infinity. For \(i:\mathbb A^1_{Y'}\hookrightarrow\mathbb P^1_{Y'}\), the Leray spectral sequence
\[
H^p(\mathbb P^1,R^qi_*F)\Rightarrow H^{p+q}(\mathbb A^1,F)
\]
reduces high-degree vanishing to the fibre of \(R^qi_*F\) at the point at infinity. By VIII, 5.2 this fibre is the cohomology of a punctured strict localization; by Assertion 4.3(d), it vanishes in degrees \(>d\). Proper base change for \(\mathbb P^1\) handles the remaining terms.
\end{prooftext}

\begin{statement}{Lemma 4.5}
Corollary 3.4(d') for all \(d'<d\) implies Assertion 4.3(d') for all \(d'\le d\).
\end{statement}

\begin{prooftext}
Let \(X'\), \(f\), and \(U\) be as in Assertion 4.3(d). Choose an algebraic model \(X_0\) and write the strict localization as the limit of affine \(\etale\) neighbourhoods \(X_i\). A constructible torsion sheaf on \(U\) descends to some \(U_i\), and VII, 5.8 identifies cohomology on \(U\) with the filtered colimit of cohomologies on the \(U_i\). It is enough to kill classes after passing to a further neighbourhood. Consider the morphism \(X_0\to\mathbb A^1\) defined by \(f\) and base-change to the strict localization of \(\mathbb A^1\) at the origin. The complement of the closed fibre is an affine algebraic scheme over the generic point and has dimension \(\le d-1\). By the induction hypothesis and Hochschild-Serre for the generic residue field, whose cohomological dimension is one, this complement has cohomological dimension \(\le d\). Hence all classes in degrees \(>d\) die after this passage, proving the assertion and completing the induction.
\end{prooftext}

% ============================================================
\clearpage
\begin{center}
{\Large \textbf{SGA 4, Expos\'e XV}}\\[0.35em]
{\large \textbf{Acyclic Morphisms}}\\[0.25em]
{\normalsize M. Artin}
\end{center}

\section*{Introduction}

Let \(g:X\to Y\) be a morphism of schemes. In the first section we study conditions ensuring that \(g\) is acyclic, meaning that for every torsion sheaf \(F\) on \(Y\), the natural maps
\[
H^q(Y,F)\longrightarrow H^q(X,g^*F)
\]
are isomorphisms, and that this remains true after replacing \(Y\) by an \(\etale\) \(Y'\to Y\). A necessary condition is that the geometric fibres be acyclic, namely have trivial cohomology for constant torsion sheaves. We shall see that this condition, together with a local condition called local acyclicity, implies acyclicity.

The fundamental theorem is that a smooth morphism is locally acyclic for torsion sheaves prime to the residual characteristics. This implies the smooth base-change theorem developed in the following expos\'e. The present expos\'e is independent of Expos\'es XI--XIV; proper base change reappears from the next expos\'e onward in combination with the theorem here.

\section*{1. Global and local acyclicity}

\begin{statement}{Proposition 1.1}
For a morphism \(g:X\to Y\), the following conditions are equivalent, with the parenthetical ``respectively'' referring to the bijective version.
\begin{enumerate}[label=(\roman*)]
\item For every \(Y'\to Y\) \(\etale\) and every sheaf of sets \(F\) on \(Y'\), the adjunction map \(F\to g'_*g'^*F\) is injective (respectively bijective).
\item For every such \(Y'\) and \(F\), the map \(H^0(Y',F)\to H^0(X',g'^*F)\) is injective (respectively bijective).
\item The same assertion holds for every locally quasi-finite \(Y'\to Y\).
\item If \(g\) is quasi-compact and quasi-separated, it is enough in (ii) to test constant sheaves \(I_{Y'}\) for one fixed set \(I\) with \(\#I\ge2\).
\end{enumerate}
\end{statement}

\begin{statement}{Lemma 1.2}
Every locally quasi-finite morphism \(Y'\to Y\) is, locally on \(Y'\) for the \(\etale\) topology, finite over an \(\etale\) neighbourhood of \(Y\). More precisely, near any point of \(Y'\) there is a diagram
\[
\xymatrix@=1.75cm{
Y'_i\ar[r]\ar[d]&Y'\ar[d]\\
Y_i\ar[r]&Y,
}
\]
where the horizontal arrows are \(\etale\), the left vertical arrow is finite, and the \(Y'_i\) cover \(Y'\).
\end{statement}

\begin{prooftext}
Take a henselization of \(Y\) at the image point. The localization of the base change of \(Y'\) at the chosen point is finite over this henselian local scheme. Descend this finite open neighbourhood to a sufficiently small \(\etale\) neighbourhood of \(Y\) by the standard limit theorem of EGA IV.
\end{prooftext}

The implication from \(\etale\) to locally quasi-finite in Proposition 1.1 follows from Lemma 1.2 and finite base change (VIII, 5.6). The equivalence with the test by constant sheaves is XII, 6.5(i).

\begin{statement}{Definition 1.3}
A morphism satisfying the injective form of Proposition 1.1 is called \((-1)\)-acyclic; one satisfying the bijective form is called \(0\)-acyclic. It is universally \((-1)\)-acyclic or universally \(0\)-acyclic if every base change has the same property. A scheme is \((-1)\)-acyclic if it is nonempty, and \(0\)-acyclic if it is connected and nonempty.
\end{statement}

\begin{statement}{Corollary 1.4}
If \(g:X\to Y\) is \(0\)-acyclic, then for every locally quasi-finite \(Y'\to Y\) and every sheaf of groups \(F\) on \(Y'\), the map
\[
H^1(Y',F)\longrightarrow H^1(X',g'^*F)
\]
is injective.
\end{statement}

\begin{statement}{Corollary 1.5}
If \(g\) is surjective, then \(g\) is \((-1)\)-acyclic. Conversely, if \(g\) is quasi-compact and quasi-separated and \((-1)\)-acyclic, then \(g\) is surjective. In that case \((-1)\)-acyclicity is already universal.
\end{statement}

\begin{statement}{Proposition 1.6}
Let \(\LL\) be a set of primes and let \(n\in\mathbb N\) (respectively \(n=1\) in the non-abelian case). The following are equivalent, with the non-abelian variant obtained by replacing abelian \(\LL\)-torsion sheaves by ind-\(\LL\)-finite groups and by restricting to degrees \(0,1\).
\begin{enumerate}[label=(\roman*)]
\item For every \(\etale\) \(Y'\to Y\) and every such sheaf \(F\) on \(Y'\),
\[
F\xrightarrow{\sim} g'_*g'^*F,
\qquad
R^qg'_*g'^*F=0\quad(1\le q\le n).
\]
\item For every \(\etale\) \(Y'\to Y\), the maps
\[
H^q(Y',F)\to H^q(X',g'^*F)
\]
are bijective for \(q\le n\) and injective for \(q=n+1\).
\item The same is true for every locally quasi-finite \(Y'\to Y\).
\item If \(g\) is quasi-compact and quasi-separated, it is enough to require surjectivity after every locally quasi-finite base change for the constant sheaves \(\Zz/\ell^\nu\Zz\), \(\ell\in\LL\), \(\nu>0\) (respectively for finite \(\LL\)-groups in \(H^0,H^1\)).
\end{enumerate}
\end{statement}

\begin{prooftext}
The equivalence of (i) and (ii) is the Leray spectral sequence, and in degree one for groups the usual non-abelian exact sequence. The passage from \(\etale\) to locally quasi-finite base changes uses Lemma 1.2 and the finite base-change isomorphism. Finally, (iv) is equivalent to (iii) by XII, 6.5.
\end{prooftext}

\begin{statement}{Definition 1.7}
A morphism satisfying Proposition 1.6 is called \(n\)-acyclic for \(\LL\) (respectively \(1\)-aspherical for \(\LL\)). It is acyclic for \(\LL\) if it is \(n\)-acyclic for all \(n\). The universal versions are defined by requiring the same property after every base change. A scheme is \(n\)-acyclic for \(\LL\) if it is \(0\)-acyclic and has \(H^q(X,\Zz/\ell\Zz)=0\) for all \(\ell\in\LL\), \(1\le q\le n\); it is \(1\)-aspherical if \(H^1(X,G)=0\) for every finite \(\LL\)-group \(G\).
\end{statement}

\begin{remarktext}{Remarks 1.8}
If \(g\) is quasi-compact and quasi-separated, universal \(n\)-acyclicity can be checked after base changes locally of finite presentation. It is not known here whether an \(n\)-acyclic morphism is automatically universally \(n\)-acyclic. For \(n=0\), Definition 1.7 recovers Definition 1.3.
\end{remarktext}

\begin{statement}{Proposition 1.9}
\begin{enumerate}[label=(\roman*)]
\item If \(X\xrightarrow gY\xrightarrow hZ\) and \(g\) is \(0\)-acyclic (respectively \(1\)-aspherical, respectively \(n\)-acyclic), then for every sheaf \(F\) on \(Z\) of the relevant type, the direct images for \(hg\) and for \(h\) are canonically identified through the relevant degrees. Consequently \(h\) has the same acyclicity property if and only if \(hg\) has it.
\item A projective limit of quasi-compact quasi-separated morphisms with affine transition maps inherits \(0\)-acyclicity, \(1\)-asphericity, and \(n\)-acyclicity from the approximating morphisms.
\end{enumerate}
\end{statement}

The proof uses the Leray spectral sequence for (i), and the limit formalism of VII, 5, for (ii).

\begin{statement}{Corollary 1.9.1}
If \(g:X\to Y\) is quasi-compact and quasi-separated and is \(n\)-acyclic (respectively \(1\)-aspherical), then every geometric fibre \(X_{\bar y}\), for \(\bar y\) algebraic over a point of \(Y\), has the same property.
\end{statement}

\begin{statement}{Proposition 1.10}
For a morphism \(g:X\to Y\), the following conditions are equivalent.
\begin{enumerate}[label=(\roman*)]
\item For every pair of geometric points \(\bar x\mapsto\bar y\), the induced morphism of strict localizations \(\widetilde X\to\widetilde Y\) is \(0\)-acyclic (respectively \(1\)-aspherical, respectively \(n\)-acyclic).
\item In every cartesian diagram
\[
\xymatrix@C=2cm@R=1.2cm{
X\ar[d]^g&X'\ar[l]\ar[d]^{g'}&X''\ar[l]\ar[d]^{g''}\\
Y&Y'\ar[l]&Y''\ar[l]
}
\]
with \(Y'\to Y\) \(\etale\) and \(Y''\to Y'\) \(\etale\) of finite presentation, the corresponding base-change maps
\[
g'^*R^qi'_*F\longrightarrow R^qj'_*g''{}^*F
\]
are isomorphisms in the prescribed range and monomorphisms in the next degree.
\item The same condition holds when \(Y''\to Y'\) is only quasi-finite and quasi-separated.
\item Equivalently, one may take \(Y'\to Y\) locally quasi-finite and \(Y''\to Y'\) a quasi-compact open immersion.
\end{enumerate}
\end{statement}

\begin{prooftext}
For (i) \(\Rightarrow\) (iii), check the base-change map on geometric fibres, where VIII, 5.2 and 5.3 identify those fibres with cohomology of strict localizations; then use the acyclicity of \(\widetilde X\to\widetilde Y\). For (ii) \(\Rightarrow\) (i), express a strict localization as the inverse limit of affine \(\etale\) neighbourhoods and descend the given \(\etale\) morphism and sheaf to one of them. The equalities of fibres furnished by VIII, 5.2 identify the desired cohomology maps. Finally, the passage between (ii), (iii), and (iv) uses the main theorem on quasi-finite morphisms: factor a quasi-finite separated morphism as an open immersion followed by a finite morphism; the finite part is handled by VIII, 5.5 and 5.8.
\end{prooftext}

\begin{statement}{Definition 1.11}
A morphism satisfying Proposition 1.10 is called locally \(0\)-acyclic, locally \(1\)-aspherical for \(\LL\), or locally \(n\)-acyclic for \(\LL\). The evident universal versions are obtained by imposing the same condition after arbitrary base change.
\end{statement}

\begin{statement}{Corollary 1.12}
If \(g:X\to Y\) is locally of finite type, it is enough to test local acyclicity at geometric points \(\bar x\) which are closed in their geometric fibre.
\end{statement}

\begin{statement}{Proposition 1.13}
\begin{enumerate}[label=(\roman*)]
\item If \(X\xrightarrow gY\xrightarrow hZ\), and \(g\) is surjective and locally \(0\)-acyclic (respectively locally \(1\)-aspherical, respectively locally \(n\)-acyclic), then \(h\) has the corresponding local property if and only if \(hg\) has it.
\item Projective limits with affine transition maps preserve local \(0\)-acyclicity, local \(1\)-asphericity, and local \(n\)-acyclicity.
\end{enumerate}
\end{statement}

\begin{remarktext}{Remarks 1.14}
The case \(n=-1\) is parallel: local \((-1)\)-acyclicity means that the induced maps on strict localizations are surjective, or equivalently that the base-change map for \(i'_*(\cdot)\) is a monomorphism. Flat morphisms are locally \((-1)\)-acyclic. For morphisms locally of finite presentation over locally noetherian bases, this is closely related to universal openness.
\end{remarktext}

\begin{statement}{Theorem 1.15}
Let \(g:X\to Y\) be quasi-compact and quasi-separated, let \(\LL\) be a set of primes, and let \(n\in\mathbb N\) (respectively \(n=1\) in the non-abelian case). Assume \(g\) is \((n-1)\)-acyclic and locally \((n-1)\)-acyclic for \(\LL\). Let \(F\) be a sheaf on \(Y\), abelian \(\LL\)-torsion if \(n\ge1\) (respectively a sheaf of groups), and let \(\xi\in H^n(X,g^*F)\). Then \(\xi\) comes from a unique element of \(H^n(Y,F)\) if and only if, for every geometric point \(\bar y\) of \(Y\) algebraic over a point of \(Y\), its restriction to \(H^n(X_{\bar y},F|_{X_{\bar y}})\) lies in the image of \(H^n(\bar y,F_{\bar y})\). For \(n\ge1\), this simply says that all those fibre restrictions are zero.
\end{statement}

\begin{statement}{Corollary 1.16}
If \(g\) is quasi-compact and quasi-separated and locally \((n-1)\)-acyclic for \(\LL\), then \(g\) is \(n\)-acyclic for \(\LL\) (respectively \(1\)-aspherical) if and only if all its geometric fibres are \(n\)-acyclic (respectively \(1\)-aspherical).
\end{statement}

\begin{statement}{Corollary 1.17}
A morphism \(g:X\to Y\) is locally \(n\)-acyclic for \(\LL\) (respectively locally \(1\)-aspherical) if and only if, for every geometric point \(\bar x\) of \(X\), all geometric fibres of the induced morphism of strict localizations \(\widetilde X\to\widetilde Y\) are \(n\)-acyclic (respectively \(1\)-aspherical).
\end{statement}

\begin{statement}{Corollary 1.18}
If \(g\) is locally of finite type, the criterion of Corollary 1.17 may be checked only at geometric points \(\bar x\) closed in their fibres.
\end{statement}

\begin{prooftext}[Proof of Theorem 1.15]
First reduce to the case \(Y\) strictly local. In the abelian case with \(n\ge1\), the Leray spectral sequence and the \((n-1)\)-acyclicity of \(g\) give an exact sequence
\[
0\to H^n(Y,F)\to H^n(X,g^*F)\to H^0(Y,R^ng_*g^*F).
\]
The last term can be tested on geometric fibres by VIII, 5.2. The case \(n=0\) and the non-abelian case \(n=1\) use respectively the same fibrewise argument for sections and the non-abelian exact sequence
\[
H^1(Y,F)\to H^1(X,g^*F)\to H^0(Y,R^1g_*g^*F).
\]

Assume henceforth \(Y\) strictly local. Corollary XII, 6.10 reduces us to finding, for each maximal point of \(Y\), a finite morphism \(Y'\to Y\) whose image contains a nonempty open neighbourhood and such that the pullback of \(\xi\) descends. Passing to a separable closure of the residue field at the maximal point provides such finite integral neighbourhoods; the fibrewise hypothesis ensures that, after one of them, the class becomes zero over the corresponding open part (or comes from sections in degree \(0\)). The complementary closed part is handled by the minimality argument of XII, 6.10.

For \(n\ge1\), after replacing \(Y\) by such a finite cover, we may suppose that \(Y=U\cup Z\) with \(\xi|_{X_U}=\xi|_{X_Z}=0\). Let \(i:U\hookrightarrow Y\), \(j:Z\hookrightarrow Y\), and use
\[
0\to F_{U,Y}\to F\to F_{Z,Y}\to0.
\]
The class \(\xi\) lifts to cohomology with coefficients in \(F_{U,Y}\), and after embedding \(F_{U,Y}\) into \(i_*F_U\) and applying XII, 6.6, it is enough to prove that
\[
H^n(X,i'_*F_{U'}')\to H^n(U',F_{U'}')
\]
is injective. The Leray spectral sequence for \(i'\) and local \((n-1)\)-acyclicity identify \(R^qi'_*F_{U'}'\) with pullbacks of \(R^qi_*F_U\) for \(q\le n-1\). Global \((n-1)\)-acyclicity then makes the relevant low-degree terms vanish, giving the desired injectivity. For \(n=1\) with non-abelian coefficients, the same argument uses the non-abelian exact sequence. For \(n=0\), glue sections over \(U\) and \(Z\); compatibility can be checked after the surjective base change by \(g\), and the local \((-1)\)-acyclicity gives the necessary injectivity of the base-change map for \(i_*\). \qedtext
\end{prooftext}

\begin{remarktext}{Remark 1.18}
There is a useful refinement when finitely many geometric points \(\bar y_i\to Y\) detect the sheaf \(F\), for example when \(Y\) is noetherian and \(F\) is constructible. In that case the criterion of Theorem 1.15 need only be checked at those points. Indeed, by XII, 6.6 one reduces to sheaves of the form \(f_*M_{\bar y}\), and local acyclicity identifies their pullbacks and higher direct images with the corresponding sheaves on the fibre \(X_{\bar y}\). The cohomological assertion then becomes precisely the assertion on that fibre.
\end{remarktext}

\section*{2. Local acyclicity of a smooth morphism}

\begin{statement}{Theorem 2.1}
Let \(g:X\to Y\) be smooth, and let \(\LL\) be the set of primes different from all residual characteristics of \(X\). Then \(g\) is locally acyclic and locally \(1\)-aspherical for \(\LL\).
\end{statement}

\begin{statement}{Corollary 2.2}
For the structural morphism \(\mathbb A^m_Y\to Y\), the morphism is acyclic and \(1\)-aspherical for the set of primes distinct from the residual characteristics of \(Y\).
\end{statement}

\begin{prooftext}
By induction on \(m\), reduce to \(m=1\). Theorem 2.1 gives the local condition; Corollary 1.16 reduces global acyclicity to the fibres. For a separably closed field \(k\), \(\Spec k[t]\) has the required prime-to-\(p\) cohomological acyclicity by IX, 4.6. Prime-to-\(p\) \(1\)-asphericity follows from the elementary fact that a finite tame Galois covering of \(\mathbb P^1_k\) unramified away from infinity is trivial: the Hurwitz formula gives
\[
2(g'-1)=2n(g-1)+\deg\mathfrak d_{X'/X}\le -2n+(n-1),
\]
forcing \(n=1\).
\end{prooftext}

\begin{prooftext}[Proof of Theorem 2.1]
The assertion is \(\etale\)-local on \(X\) and \(Y\). We may therefore take \(X=\mathbb A^m_Y\), and then reduce by transitivity to \(m=1\): \(X=\Spec\Oo_Y[t]\). By Corollary 1.18 it is enough to check the fibres of the induced morphism between strict localizations at a point \(\bar x\) closed in its geometric fibre. After a finite local base change on \(Y\), one may assume the residue fields of \(\bar x\) and \(\bar y\) coincide. Translating the coordinate, one may assume \(\bar x\) is the point \(t=0\). Thus one is reduced to rings of the form \(A\{t\}\), the henselization of \(A[t]\) at the ideal \((\operatorname{rad}A,t)\).
\end{prooftext}

\begin{statement}{Notation 2.3}
For a henselian local ring \(A\), denote by \(A\{t\}\) the henselization of \(A[t]\) at the prime ideal generated by \(\operatorname{rad}A\) and \(t\).
\end{statement}

\begin{statement}{Corollary 2.4}
If \(A\) is henselian and \(A'\) is a finite \(A\)-algebra, then
\[
A'\{t\}\simeq A'\otimes_A A\{t\}.
\]
\end{statement}

\begin{statement}{Lemma 2.5}
Let \(A\) be henselian with residual characteristic \(p\). Every geometric fibre of \(\Spec A\{t\}\to\Spec A\) is acyclic and \(1\)-aspherical for \(\Primes-\{p\}\).
\end{statement}

\begin{prooftext}
Each geometric fibre is a projective limit of affine algebraic curves over a separably closed field. Hence its cohomological dimension is at most one by IX, 5.7 and VII, 5.7. It remains to prove \(1\)-asphericity. Write \(A\) as a filtered limit of henselizations of finite-type \(\Zz\)-algebras at prime ideals. Since \(A\{t\}=\varinjlim A_\alpha\{t\}\), limit arguments reduce to the case where \(A\) is excellent. Nonemptiness follows from the section \(t=0\). Connectedness is reduced, after normalization in finite separable field extensions and using Corollary 2.4, to the fact that \(A\{t\}\) is integral when \(A\) is normal. The remaining assertion is Lemma 2.6.
\end{prooftext}

\begin{statement}{Lemma 2.6}
Let \(A\) be henselian and excellent, and let \(\bar Z\) be a geometric fibre of \(\Spec A\{t\}\to\Spec A\). Every principal Galois cover of \(\bar Z\) with group of order prime to the residual characteristic of \(A\) is trivial.
\end{statement}

\section*{3. Proof of the principal lemma}

After replacing \(A\) by a finite normal algebra, Lemma 2.6 reduces to the generic fibre of \(\Spec A\{t\}\), with \(A\) normal. It is enough to prove the following form.

\begin{statement}{Lemma 3.1}
Let \(A\) be normal henselian excellent with residual characteristic \(p\), and let \(Z\) be the generic fibre of \(\Spec A\{t\}\to\Spec A\). If \(Z'\to Z\) is a principal finite Galois \(\etale\) covering of group order prime to \(p\), then \(Z'\to Z\) is trivial.
\end{statement}

\begin{prooftext}
Let \(K\) be the fraction field of \(A\). After finite separable extension of \(K\) and normalization of \(A\), the cover descends to a finite normal \(A\{t\}\)-algebra. The ramification is prime to the residual characteristic and is controlled by the divisor \(t=0\). The argument reduces, by normalization and localization, to the following purity-style lemma.
\end{prooftext}

\begin{statement}{Lemma 3.2}
Let \(B\) be a normal excellent local ring, \(0\ne t\in\operatorname{rad}B\), and let \(B'\) be a finite normal integral \(B\)-algebra containing \(B\). Assume the induced extension is \(\etale\) away from the closed point and splits completely on \(\Spec(B/tB)\) away from the closed point. Then, after localization at the possible non-\(\etale\) locus, one is reduced to the case \(\dim B\ge3\), where Lemma 3.3 applies.
\end{statement}

\begin{statement}{Lemma 3.3}
Let \(B\) be an excellent normal local ring, let \(0\ne t\in\operatorname{rad}B\), and assume \(\dim B\ge3\). Let \(B'\) be a finite normal integral \(B\)-algebra containing \(B\), and assume \(B'/tB'\) is completely split over \(B/tB\) away from the closed point. Then \(B'\) is \(\etale\) over \(B\).
\end{statement}

\begin{prooftext}
Let \(X=\Spec B\) and let \(U\subset X\) be the open where \(B'\) is \(\etale\). Since \((X-U)\cap V(t)\) is just the closed point, and \(V(t)\) is defined by one equation, one has \(\dim(X-U)\le1\). Because \(X\) is catenary and \(\dim B\ge3\), the complement has codimension at least two. For \(X'=\Spec B'\), universal catenarity gives the same codimension estimate for the complement of \(U'=X'\times_XU\), and normality gives
\[
B'=\Gamma(X',\Oo_{X'})\simeq\Gamma(U',\Oo_{U'}).
\]
Equivalently, the finite algebra is the direct image of its \(\etale\) restriction to \(U\). Complete \(B\). The completion is normal and faithfully flat, and direct image from \(U\) commutes with this base change. Thus it is enough to treat the complete case. In that case the local Lefschetz theorem of SGA 2, X, 2.1 and 2.5, applied with the depth estimate just obtained, shows that any such cover of \(U\) split on the punctured divisor is itself completely split. Hence it extends as a trivial \(\etale\) covering over \(X\), so \(B'\) is \(\etale\) over \(B\).
\end{prooftext}

\begin{remarktext}{Remark 3.4}
The proof of Lemmas 3.2 and 3.3 works under more general hypotheses. It is enough, for example, that the completion \(\widehat B\) be normal, or even that the \((tB)\)-adic completion be normal and a quotient of a regular ring.
\end{remarktext}

\section*{4. Appendix: a criterion for local \(0\)-acyclicity}

The results of this section are not used later in the seminar.

\begin{statement}{Theorem 4.1}
Let \(g:X\to Y\) be flat with separable fibres. Assume \(X\) and \(Y\) locally noetherian (respectively assume \(g\) locally of finite presentation). Then \(g\) is universally locally \(0\)-acyclic.
\end{statement}

\begin{prooftext}
By standard limit arguments, reduce to the locally noetherian case and then to proving local \(0\)-acyclicity. Corollary 1.17 reduces this to the assertion that the geometric fibres of a strictly local flat morphism with separable fibres are connected. This is EGA IV, 18.9.8.
\end{prooftext}

\begin{statement}{Corollary 4.2}
If \(Y\) is discrete, then every morphism \(g:X\to Y\) is universally locally \(0\)-acyclic.
\end{statement}

Reduce to \(Y=\Spec k\) with \(k\) perfect, then approximate \(X\) by schemes of finite type over \(k\), replace by its reduction, and apply Theorem 4.1.

\begin{statement}{Corollary 4.3}
Suppose \(g\) is surjective and satisfies the hypotheses of Theorem 4.1 or Corollary 4.2; in the first, non-finite-presentation form, assume also that \(g\) is either quasi-compact and quasi-separated or universally open. Then \(g\) is a universal effective descent morphism for the fibred category of \(\etale\) sheaves on variable schemes.
\end{statement}

\begin{prooftext}
Such a morphism is universally submersive, so VIII, 9.1 gives universal descent. Effectivity is local on the target, and after replacing \(X\) by a finite disjoint union of affine opens whose images cover the affine target, one may suppose \(X\) affine and quasi-compact over \(Y\). The argument of VIII, 9.4.1 applies, using the smooth-base-change result that \(g_*F\) commutes with locally \(0\)-acyclic base change.
\end{prooftext}

\begin{remarktext}{Remarks 4.4}
This gives the deferred proof of VIII, 9.4(d). A simpler proof in the case of field extensions reduces to \(\Spec K\to\Spec k\), for which universal local \(0\)-acyclicity follows from the prime-to-characteristic acyclicity discussed later. It is plausible that every quasi-compact, quasi-separated, surjective, locally \((-1)\)-acyclic morphism is an effective descent morphism for \(\etale\) sheaves. It is also plausible that under the hypotheses of Corollary 4.2 one has local acyclicity and local \(1\)-asphericity for all primes distinct from the residual characteristics; in characteristic zero this follows from resolution of singularities.
\end{remarktext}



\clearpage
\begin{center}
{\Large \textbf{SGA 4, Expos\'e XVI}}\\[0.35em]
{\large \textbf{The Smooth Base-Change Theorem and Applications}}\\[0.25em]
{\normalsize M. Artin}
\end{center}

\section{The smooth base-change theorem}

Let
\[
\begin{array}{ccc}
X'&\xrightarrow{g'}&X\\
\downarrow f'&&\downarrow f\\
S'&\xrightarrow{g}&S
\end{array}
\]
be a cartesian square of schemes. The base-change morphism
\[
g^*R^qf_*F\longrightarrow R^qf'_*g'^*F
\]
has already been constructed in Expos\'e XII. The purpose of the present section is to prove that it is an isomorphism when the lower morphism \(g\) is smooth and the coefficients are torsion prime to the residual characteristics.

\begin{statement}{Theorem 1.1}
Let \(f:X\to S\) be any morphism, let \(g:S'\to S\) be smooth, and let \(F\) be a torsion abelian sheaf on \(X\) whose torsion primes are invertible on \(S'\). Then the canonical base-change homomorphism
\[
g^*R^qf_*F\longrightarrow R^qf'_*g'^*F
\]
is an isomorphism for every \(q\). The corresponding assertions in degree zero for sheaves of sets and in degree one for sheaves of groups hold in the same prime-to-characteristic range.
\end{statement}

\begin{prooftext}
The assertion is local on \(S'\) for the etale topology. Since \(g\) is smooth, etale-locally it factors as an etale morphism followed by a projection \(\mathbb A^r_S\to S\). Etale base change is formal. Thus the essential case is the projection from affine space. By Expos\'e XV, a smooth morphism is locally acyclic for torsion primes invertible on the base. Applying local acyclicity to the strict localizations used to compute the stalks of \(R^qf_*F\), one obtains the required isomorphism on every geometric fibre of \(S'\). Since the assertion may be checked on stalks, the theorem follows. \qedtext
\end{prooftext}

\begin{statement}{Corollary 1.2}
If \(S'\to S\) is etale, then the same base-change isomorphism holds for all torsion coefficients. In particular, higher direct images for the etale topology are compatible with etale localization on the base.
\end{statement}

\begin{statement}{Corollary 1.3}
If \(S'\to S\) is a localization, a strict localization, or the inclusion of a geometric point obtained as a limit of smooth neighborhoods in the prime-to-characteristic range, the same comparison remains valid by passage to the limit.
\end{statement}

\begin{remarktext}{Remark 1.4}
The theorem is not a substitute for proper base change. Proper base change allows arbitrary change of base for proper \(f\); smooth base change allows a special class of base changes for arbitrary \(f\). The two results are complementary and are combined repeatedly in later arguments.
\end{remarktext}

\section{Specialization theorem for cohomology groups}

Let \(S\) be a scheme, let \(s\leadsto t\) be a specialization represented by a morphism from a strictly local scheme, and let \(X\to S\) be a morphism. Smooth base change gives canonical specialization maps between the cohomology groups of the geometric fibres, whenever the specialization is realized through smooth neighborhoods and the coefficients are prime to the residual characteristics.

\begin{statement}{Theorem 2.1}
Let \(f:X\to S\) be a morphism and let \(F\) be a torsion sheaf on \(X\), with torsion primes invertible on \(S\). Suppose that \(S\) is locally noetherian and that two geometric points \(\bar s\) and \(\bar t\) lie in a common smooth specialization situation. Then there are canonical specialization homomorphisms
\[
H^q(X_{\bar s},F_{\bar s})\longrightarrow H^q(X_{\bar t},F_{\bar t})
\]
functorial in \(F\), compatible with cup products and long exact sequences, and invariant under refinement of the chosen smooth neighborhood.
\end{statement}

The proof is a direct reading of the stalks of \(R^qf_*F\). Smooth base change identifies the pullback of \(R^qf_*F\) to a smooth neighborhood with the higher direct image after pullback. The strict local base has a closed point and a generization; evaluation at these points gives the specialization map.

\begin{statement}{Corollary 2.2}
If \(f:X\to S\) is smooth and proper, and \(F\) is finite locally constant with order invertible on \(S\), then \(R^qf_*F\) is locally constant finite. Hence the cohomology groups of the geometric fibres form a locally constant system on \(S\).
\end{statement}

Proper base change gives constructibility and fibrewise identification; smooth base change shows that the constructible sheaf has no specialization jumps on the smooth base. After shrinking to connected components, the finite stalks and specialization maps are constant.

\begin{statement}{Corollary 2.3}
For a smooth proper family over a connected normal base, the prime-to-characteristic cohomology groups of the geometric fibres are canonically identified up to the monodromy action of the etale fundamental group of the base.
\end{statement}

This is the etale analogue of the topological fact that a proper smooth map of complex manifolds is locally topologically trivial.

\section{Relative cohomological purity}

The next application is a relative purity statement. Let \(X\) be regular, let \(Y\subset X\) be a regular closed subscheme of codimension \(c\), and let \(U=X\setminus Y\). For torsion coefficients prime to the residual characteristics, the local cohomology of \(X\) with support in \(Y\) is concentrated in degree \(2c\) and is given by a Tate twist.

\begin{statement}{Theorem 3.1}
Let \(i:Y\hookrightarrow X\) be a regular immersion of codimension \(c\), with \(X\) and \(Y\) locally noetherian, and let \(n\) be invertible on \(X\). Then
\[
R^qi^!\Zz/n\Zz=0\quad(q\ne 2c),
\qquad
R^{2c}i^!\Zz/n\Zz\simeq \Zz/n\Zz(-c).
\]
Equivalently,
\[
H^q_Y(X,\Zz/n\Zz)=0\quad(q<2c),
\]
and the first non-zero local cohomology sheaf is the expected orientation sheaf.
\end{statement}

\begin{prooftext}
The assertion is local for the etale topology on \(X\). By the definition of a regular immersion, and after etale localization, the pair \((X,Y)\) is compared to \((\mathbb A^c_Y,Y)\), where \(Y\) is embedded as the zero section. Smooth base change reduces the calculation to the case where \(Y\) is the spectrum of a separably closed field. Then the computation is the cohomology with supports of affine space at the origin. For \(c=1\) this is the Kummer sequence and the calculation of \(\mathbb G_m\); the general case follows by iterating the one-dimensional calculation and using the Kunneth formula for supports. \qedtext
\end{prooftext}

\begin{statement}{Corollary 3.2}
If \(X\) is regular and \(U=X\setminus Y\), with \(Y\) regular of codimension \(c\), then the first possible obstruction to extending a prime-to-characteristic torsion cohomology class from \(U\) across \(Y\) occurs in degree \(2c-1\). In codimension at least two, low-degree cohomology is unchanged by removing \(Y\).
\end{statement}

\begin{remarktext}{Remark 3.3}
In later language this is absolute cohomological purity. In this expos\'e it is obtained in the relative form needed for base-change and finiteness applications.
\end{remarktext}

\section{Comparison theorem over the complex numbers}

Let \(X\) be a scheme locally of finite type over \(\Cc\), and let \(X^{\an}\) be the associated complex analytic space. A sheaf \(F\) on \(X_{\et}\) gives, by pullback along the natural functor from analytic neighborhoods to etale neighborhoods, a sheaf \(F^{\an}\) on \(X^{\an}\).

\begin{statement}{Theorem 4.1}
For every constructible torsion sheaf \(F\) on \(X_{\et}\), the natural comparison maps
\[
H^q(X_{\et},F)\longrightarrow H^q(X^{\an},F^{\an})
\]
are isomorphisms for all \(q\). The same is true for cohomology with supports and is compatible with long exact sequences, cup products, and proper or smooth base change.
\end{statement}

The smooth case with locally constant finite coefficients was treated in Expos\'e XI. The general case follows by embedding \(X\) into a smooth scheme locally, stratifying so that \(F\) is locally constant finite on the strata, and using the exact sequence for an open stratum and its closed complement. Proper base change and smooth base change ensure compatibility with the two geometric operations needed in the induction.

\begin{statement}{Corollary 4.2}
The etale cohomological dimension bounds proved algebraically agree over \(\Cc\) with the corresponding classical bounds for the analytic topology.
\end{statement}

\section{Finiteness theorem for algebraic schemes over a separably closed field}

\begin{statement}{Theorem 5.1}
Let \(X\) be a scheme of finite type over a separably closed field \(k\), and let \(F\) be a constructible torsion sheaf on \(X\). Then the groups \(H^q(X,F)\) are finite when \(F\) has finite fibres, and they vanish for all sufficiently large \(q\). More precisely, the vanishing is bounded by the cohomological-dimension estimates of Expos\'es X and XIV.
\end{statement}

\begin{prooftext}
Choose a compactification of \(X\) when available, or argue by noetherian induction using an open affine stratification. On affine strata the affine Lefschetz theorem gives the dimension bound. On proper pieces, the finiteness theorem for proper morphisms applies. The long exact sequence for an open and its closed complement then gives both finiteness and vanishing for \(X\). For a non-algebraically closed but separably closed field, the same proof applies after using the invariance of the etale topos under purely inseparable extension. \qedtext
\end{prooftext}

\begin{statement}{Corollary 5.2}
Let \(f:X\to S\) be a morphism of finite type with \(S\) noetherian and let \(F\) be constructible torsion on \(X\). After stratifying \(S\), the sheaves \(R^qf_*F\) are constructible in the range controlled by the preceding theorem whenever \(f\) is proper, smooth, or is built from these cases by the standard open-closed decompositions.
\end{statement}

This corollary summarizes the role of Expos\'e XVI: smooth base change supplies local constancy in smooth directions, proper base change supplies invariance under proper specialization, and the affine cohomological-dimension theorem supplies vanishing on open affine pieces.

\clearpage
\begin{center}
{\Large \textbf{SGA 4, Expos\'e XVII -- Beginning}}\\[0.35em]
{\large \textbf{Cohomology with Proper Supports}}\\[0.25em]
{\normalsize P. Deligne}
\end{center}

\section*{Introduction}
\addcontentsline{toc}{section}{XVII. Introduction}

The preceding expos\'es developed the basic finiteness, base-change, and acyclicity theorems for torsion etale cohomology. The present expos\'e begins the construction of cohomology with proper supports and of the functors which will ultimately be denoted \(Rf_!\). The construction is not merely a formal variant of ordinary direct image. It requires a derived-category formalism capable of handling compactifications, gluing, and independence of choices.

The guiding principle is the following. If \(f:X\to S\) is separated and of finite type, choose a compactification
\[
X\xrightarrow{j}\bar X\xrightarrow{\bar f}S,
\]
with \(j\) an open immersion and \(\bar f\) proper. Then cohomology with proper supports should be computed by
\[
Rf_!K=R\bar f_*j_!K.
\]
One must prove that this object is independent of the compactification, functorial in \(f\), compatible with base change, and compatible with the operations needed in duality. Expos\'e XVII supplies the formal machinery for this construction.

\section{Terminological preliminaries}

\subsection{Categories, universes, and topoi}
All categories and topoi are understood relative to the fixed universe conventions of the earlier expos\'es. A topos is a category equivalent to the category of sheaves on a site. Morphisms of topoi are written with the usual pair of adjoint functors \((f^*,f_*)\), where \(f^*\) is left exact. For a ringed topos \((X,A)\), modules always mean sheaves of \(A\)-modules on \(X\).

\subsection{Complexes}
A complex \(K\) in an abelian category \(\A\) is cohomologically indexed:
\[
\cdots\to K^{n-1}\xrightarrow{d^{n-1}}K^n\xrightarrow{d^n}K^{n+1}\to\cdots,
\qquad d^{n}d^{n-1}=0.
\]
The translated complex \(K[1]\) is defined by \(K[1]^n=K^{n+1}\) with differential \(-d_K\). This sign is essential for the standard triangles and mapping cones to behave functorially.

\subsection{Homotopy and derived categories}
Let \(C(\A)\) be the category of complexes, \(K(\A)\) the homotopy category, and \(D(\A)\) the derived category obtained by inverting quasi-isomorphisms. The bounded variants are denoted \(D^+(\A)\), \(D^-(\A)\), and \(D^b(\A)\). A functor between abelian categories which is left exact, right exact, or exact gives rise, under the usual hypotheses on injective or flat resolutions, to derived functors between the corresponding derived categories.

\subsection{Triangulated conventions}
Distinguished triangles are written
\[
K\longrightarrow L\longrightarrow M\longrightarrow K[1].
\]
A short exact sequence of complexes gives such a triangle. Conversely, the cohomology long exact sequence attached to a distinguished triangle is the principal computational device. The octahedral axiom is used only in its standard consequence: mapping cones and compositions are compatible up to unique isomorphism in the derived category.

\subsection{Truncations}
For a complex \(K\), the canonical truncations \(\tau_{\le n}K\) and \(\tau_{>n}K\) fit into a distinguished triangle
\[
\tau_{\le n}K\to K\to \tau_{>n}K\to \tau_{\le n}K[1].
\]
They give the standard spectral sequences used throughout the expos\'e. In particular, if a cohomological functor is known on sheaves placed in one degree, the truncation filtration reduces statements about bounded below complexes to statements about their cohomology sheaves.

\subsection{Flat and injective resolutions}
In the categories of sheaves of modules on a ringed topos, there are enough injectives. Flat resolutions exist under the hypotheses used here and may be chosen functorially enough for derived inverse image and tensor product. When a functor is exact on injectives, or sends a chosen class of adapted objects to adapted objects, its derived functor is computed by resolving only once.

\subsection{Notation for supports}
If \(j:U\hookrightarrow X\) is an open immersion and \(i:Z\hookrightarrow X\) is the closed complement, \(j_!\) denotes extension by the initial object outside \(U\) for sheaves of sets, and extension by zero for abelian sheaves or modules. The fundamental exact sequence is
\[
0\to j_!j^*F\to F\to i_*i^*F
\]
with the usual interpretation in the abelian case; it gives the localization triangle
\[
Rj_!j^*K\to K\to Ri_*i^*K\to.
\]

\subsection{Compactifiable morphisms}
A morphism \(f:X\to S\) is compactifiable if it factors as an open immersion \(j:X\hookrightarrow \bar X\) followed by a proper morphism \(\bar f:\bar X\to S\). For schemes, separated morphisms of finite type are compactifiable by Nagata's compactification theorem in the contexts in which it is used here. The construction of \(Rf_!\) first depends on the choice of such a compactification and then proves independence.

\section{Derived categories}

\subsection{Exact functors and the sign rules}

\begin{statement}{1.1.1}
Let \(F:\A\to\B\) be an additive functor between abelian categories. Applying \(F\) degreewise gives a functor on complexes. If \(F\) is exact, it sends quasi-isomorphisms to quasi-isomorphisms and therefore induces a triangulated functor
\[
D(\A)\longrightarrow D(\B).
\]
It commutes with shifts and carries distinguished triangles coming from short exact sequences of complexes to distinguished triangles.
\end{statement}

The compatibility with shifts is not literally automatic unless the sign convention for \([1]\) is fixed. With the convention \(d_{K[1]}=-d_K\), the natural isomorphism \(F(K[1])\simeq F(K)[1]\) is compatible with differentials.

\begin{statement}{*1.1.2}
For a morphism of complexes \(u:K\to L\), the cone \(C(u)\) is the complex
\[
C(u)^n=L^n\oplus K^{n+1},
\qquad
 d(l,k)=(d_Ll+u(k),-d_Kk).
\]
There is a distinguished triangle
\[
K\xrightarrow{u}L\to C(u)\to K[1].
\]
A morphism \(u\) is a quasi-isomorphism if and only if \(C(u)\) is acyclic.
\end{statement}

\begin{statement}{1.1.4}
Let \(F\) and \(G\) be exact functors. A natural transformation \(F\to G\) induces a natural transformation between the associated derived-category functors. If the transformation is an isomorphism on the abelian category, it is an isomorphism on the derived category.
\end{statement}

\begin{statement}{1.1.4.2.1}
Given two composable morphisms of complexes \(K\xrightarrow{u}L\xrightarrow{v}M\), the cones \(C(u)\), \(C(v)\), and \(C(vu)\) fit into the standard octahedron. All later compatibility statements for boundary morphisms are consequences of this octahedron.
\end{statement}

\begin{statement}{1.1.5}
If \(0\to K'\to K\to K''\to0\) is a short exact sequence of complexes, then there is a functorial distinguished triangle
\[
K'\to K\to K''\to K'[1].
\]
The connecting morphism \(K''\to K'[1]\) is represented by the usual boundary map; the sign is the one determined by the cone convention above.
\end{statement}

\begin{statement}{1.1.6}
Let \(F:\A\to\B\) be left exact. Its right derived functor \(RF:D^+(\A)\to D^+(\B)\) is constructed by replacing a bounded below complex by an injective resolution and applying \(F\). If \(F\) sends injective objects to objects acyclic for another left exact functor \(G\), then
\[
R(GF)\simeq RG\,RF.
\]
\end{statement}

\begin{statement}{1.1.8}
Dually, if \(F\) is right exact and enough flat or projective resolutions are available, one constructs \(LF\) on \(D^-(\A)\). In sheaf theory the important examples are derived tensor product and derived inverse image for modules over ringed topoi.
\end{statement}

\begin{statement}{1.1.9}
For a ringed topos \((X,A)\), the tensor product of complexes is given by the total complex of the double complex \(K^p\otimes_A L^q\). The commutativity constraint involves the Koszul sign
\[
x\otimes y\longmapsto (-1)^{pq}y\otimes x,
\]
for homogeneous elements \(x\in K^p\), \(y\in L^q\). With this convention, derived tensor product is symmetric monoidal.
\end{statement}

\begin{statement}{*1.1.10.1}
If \(u:K\to K'\) and \(v:L\to L'\) are morphisms of complexes, the tensor product of cones is compatible with the cone of the tensor product morphism only after inserting the standard Koszul signs. These signs are the source of the signs in the projection formula and in the Kunneth morphism.
\end{statement}

\begin{statement}{1.1.12}
Let \(f:(X,A)\to(Y,B)\) be a morphism of ringed topoi. The inverse image of modules is right exact and its left derived functor is denoted
\[
Lf^*:D^-(B)\to D^-(A).
\]
When the objects involved are flat, this is computed by the ordinary inverse image. The direct image \(f_*\) is left exact and has a right derived functor \(Rf_*\) on bounded below derived categories.
\end{statement}

\begin{statement}{1.1.15}
For \(K\in D^-(B)\) and \(L\in D^+(A)\) there is a natural projection morphism
\[
K\otimes_B^L Rf_*L\longrightarrow Rf_*(Lf^*K\otimes_A^L L).
\]
Under the usual boundedness and finite Tor-dimension hypotheses this morphism is an isomorphism. Establishing such isomorphisms for \(Rf_!\) is one of the purposes of the later sections.
\end{statement}

\begin{statement}{1.1.16}
All constructions above are compatible with restriction to open subtopoi, with extension by zero, and with the localization triangles associated with an open-closed decomposition. These compatibilities are fixed once and for all by the cone and shift conventions; no further sign choices will be made later.
\end{statement}


\subsection{Derived functors}

Let \(\A\) be a triangulated category and let \(S\) be a saturated multiplicative system. If \(A\in\Ob(\A)\), we denote by \(A^+\) and \(A^-\) the ind- and pro-objects
\[
 A^+=\varinjlim_{A\xrightarrow{s}A'} A',
 \qquad
 A^-=\varprojlim_{A'\xrightarrow{s}A}A',
\]
where \(s\) ranges over the filtered categories of arrows in \(S\) with source, respectively target, \(A\). The functors \(A\mapsto A^+\) and \(A\mapsto A^-\) invert the elements of \(S\), and they identify the localized category \(\A(S^{-1})\) with full subcategories of the categories of ind- and pro-objects of \(\A\). This is the device used here to define derived functors without choosing resolutions at the outset.

\begin{statement}{Definition 1.2.1}
Let \(F\) be an exact covariant multi-functor from triangulated categories \((\A_i)_{0<i\le n}\) to a triangulated category \(\A\), and let \(S_i\) and \(S\) be saturated multiplicative systems in \(\A_i\) and \(\A\).
\begin{enumerate}[label=(\roman*)]
\item The right derived functor \(RF\), relative to the \(S_i\) and \(S\), is the functor from the localized categories \(\A_i(S_i^{-1})\) to the category of ind-objects of \(\A(S^{-1})\) obtained by sending \((A_i)\) to the image of \(F(A_i^+)\).
\item The functor \(RF\) is said to be defined at \((A_i)\) if the resulting ind-object of \(\A(S^{-1})\) is essentially constant, that is, if it comes from an actual object of \(\A(S^{-1})\). That object is then called the value of \(RF\) at \((A_i)\).
\item The functor \(F\) is right derivable if \(RF\) is defined everywhere. In that case the same symbol \(RF\) denotes the resulting functor with values in \(\A(S^{-1})\).
\item A family \((A_i)\) is said to be adapted, or split, for \(RF\) if the canonical morphism \(F(A_i)\to RF(A_i)\) is an isomorphism.
\end{enumerate}
\end{statement}

The left derived functor \(LF\) is defined dually, with values first in pro-objects. The same definitions apply when \(F\) is covariant in some variables and contravariant in others. When the derived functor is everywhere defined, this definition agrees with Verdier's definition of derived functors in localized triangulated categories.

\begin{statement}{Proposition 1.2.2}
Let \((A'_1,A_1,A''_1)\) be a distinguished triangle in \(\A_1\), and let \(A_i\in\Ob(\A_i)\) for \(i\ne1\).
\begin{enumerate}[label=(\roman*)]
\item If \((A'_1,(A_i))\) and \((A''_1,(A_i))\) are split for \(RF\), then \((A_1,(A_i))\) is split for \(RF\).
\item If \(RF\) is defined at \((A'_1,(A_i))\) and at \((A''_1,(A_i))\), then it is defined at \((A_1,(A_i))\), and
\[
 RF(A'_1,(A_i))\to RF(A_1,(A_i))\to RF(A''_1,(A_i))\to RF(A'_1,(A_i))[1]
\]
is a distinguished triangle in \(\A(S^{-1})\).
\end{enumerate}
\end{statement}

The implication (ii) \(\Rightarrow\) (i) follows at once from exactness of \(F\): the comparison from the triangle obtained by applying \(F\) to the triangle obtained by applying \(RF\) is an isomorphism at two vertices, hence at the third.

\begin{statement}{Lemma 1.2.2.1}
If \(\Delta=(X,Y,Z)\) is a distinguished triangle in a triangulated category \(\A\) and \(S\) is a saturated multiplicative system, then the triangle \(\Delta^+=(X^+,Y^+,Z^+)\) in the ind-category is a filtered inductive limit of distinguished triangles of \(\A\).
\end{statement}

Let \(\mathcal L\) be the category whose objects are morphisms of distinguished triangles with source \(\Delta\) and whose three component arrows belong to \(S\). The usual calculus of fractions for saturated systems gives the following facts. If \(\Delta'\) is another distinguished triangle and \(f:\Delta\to\Delta'\) is a morphism in \(\A(S^{-1})\), then after composing \(\Delta'\) with a morphism of triangles whose components lie in \(S\), the morphism \(f\) is represented by an actual morphism of triangles in \(\A\). If two such representatives become equal in \(\A(S^{-1})\), then they become equal after another such composition. These two assertions imply that \(\mathcal L\) is filtered, and
\[
 \Delta^+=\varinjlim_{s\in\mathcal L}\operatorname{target}(s).
\]
This proves the lemma.

To prove Proposition 1.2.2(ii), choose objects \(X'\) and \(X''\) representing \(RF(A'_1,(A_i))\) and \(RF(A''_1,(A_i))\), and complete the degree-one morphism \(X''\to X'\) to a distinguished triangle \((X',X,X'')\). After replacing the input objects by targets of morphisms in the relevant multiplicative systems, the given comparison maps are represented by sections. The axiom TR3 then completes them to a morphism of distinguished triangles. Applying \(\Hom(Y,-)\) for arbitrary \(Y\in\A(S^{-1})\), one obtains a morphism between two exact sequences. The four outside vertical arrows are isomorphisms by the representing property of \(X'\) and \(X''\), hence the central arrow is an isomorphism by the five lemma. Therefore \(X\) represents \(RF(A_1,(A_i))\), and the displayed triangle is distinguished.

\begin{statement}{Definition 1.2.3}
Let \(\A\) and \(\B\) be abelian categories and let \(F:\A\to\B\) be an additive covariant functor.
\begin{enumerate}[label=(\roman*)]
\item The right derived functors of \(F\) are the functors
\[
D^+(\A)\to \operatorname{Ind}D^+(\B),
\qquad
D(\A)\to \operatorname{Ind}D(\B),
\]
derived from the extensions of \(F\) to complexes.
\item An object \(A\in\A\) is right acyclic for \(F\), or acyclic for \(RF\), if the complex with \(A\) in degree zero is split for \(RF\).
\end{enumerate}
\end{statement}

The main case is the one in which \(RF\) is everywhere defined, so that it may be regarded as a functor to \(D^+(\B)\), or to \(D(\B)\), not merely to the ind-category.

\begin{statement}{Proposition 1.2.4}
The natural square
\[
\begin{array}{ccc}
D^+(\A)&\longrightarrow&D(\A)\\
\downarrow RF&&\downarrow RF\\
\operatorname{Ind}D^+(\B)&\longrightarrow&\operatorname{Ind}D(\B)
\end{array}
\]
is commutative.
\end{statement}

Indeed, if \(K\) is zero in degrees \(<n\) and \(K\to L\) is a quasi-isomorphism, then the composite \(K\to L\to \tau_{\ge n}L\) is again a quasi-isomorphism. Thus quasi-isomorphisms from \(K\) to bounded below complexes form a cofinal subcategory of all quasi-isomorphisms with source \(K\). This proves the proposition formally. In particular, if \(RF:D(\A)\to D(\B)\) is everywhere defined, its restriction to \(D^+(\A)\) is the usual bounded-below derived functor.

\begin{statement}{1.2.5}
The preceding definitions and Proposition 1.2.4 extend without change to bounded complexes of multi-functors, covariant in some variables and contravariant in others. The dual statements define left derived functors and compare \(D^-(\A)\) with \(D(\A)\).
\end{statement}

\begin{statement}{Proposition 1.2.6}
Let \(\A\) and \(\B\) be abelian categories, let \(\mathcal P\subset\Ob(\A)\), and let \(F:\A\to\B\) be additive. Suppose that every object of \(\A\) is a quotient of an object of \(\mathcal P\), and that whenever
\[
0\to P\to Q\to R\to0
\]
is exact with \(Q,R\in\mathcal P\), one has \(P\in\mathcal P\) and
\[
0\to FP\to FQ\to FR\to0
\]
is exact. Then every object of \(\mathcal P\) is acyclic for \(LF\).
\end{statement}

The left resolutions of an object by elements of \(\mathcal P\) are cofinal among all left resolutions. If a complex
\[
\cdots\to A_2\to A_1\to A_0\to A\to0
\]
with \(A_i\in\mathcal P\) is acyclic, cut it into short exact sequences
\[
0\to\operatorname{Im}(d_{i-1})\to A_i\to\operatorname{Im}(d_i)\to0.
\]
Induction on \(i\) shows that each image belongs to \(\mathcal P\), and that applying \(F\) preserves these short exact sequences. Hence the image complex is acyclic.

\begin{statement}{Proposition 1.2.7}
Let \(F:\A\to\B\) be additive. If every object of \(\A\) is a quotient, respectively a subobject, of an object acyclic for \(LF\), respectively for \(RF\), then
\[
LF:D^-(\A)\to D^-(\B),
\qquad
RF:D^+(\A)\to D^+(\B)
\]
are everywhere defined. Moreover every bounded-above, respectively bounded-below, complex of acyclic objects is split for the corresponding derived functor.
\end{statement}

For \(LF\), quasi-isomorphisms from complexes of acyclic objects to a given bounded-above complex form a cofinal system. It remains to see that if \(s:L\to K\) is such a quasi-isomorphism between complexes of acyclic objects, then \(F(s)\) is a quasi-isomorphism. The cone of \(s\) is acyclic and still has acyclic terms. Its brutal truncations are bounded complexes of acyclic objects, hence split; passing to cohomology and letting the truncation bound tend to infinity gives the desired vanishing. The proof for \(RF\) is dual.

\begin{statement}{Proposition 1.2.8}
Let \(F^\bullet\) be a bounded-below complex of functors from \(\A\) to \(\B\). Assume that for every \(K\in K^+(\A)\) there exists a quasi-isomorphism \(L\to K\) with \(L\in K^+(\A)\) split for all the functors \(RF^i\). Then \(RF^\bullet:D^+(\A)\to D^+(\B)\) exists, and every bounded-below complex split for all the \(RF^i\) is split for \(RF^\bullet\).
\end{statement}

This is the usual spectral-sequence argument applied to the double complex obtained by resolving termwise. The dual statement holds for left derived functors and bounded-above complexes.

\begin{statement}{Definition 1.2.9}
Let \(F:D^b(\A)\to D(\B)\) be a functor and let \(d\subset\Zz\) be an interval. We say that \(F\) has cohomological amplitude contained in \(d\) if, for every object \(A\in\A\),
\[
H^iF(A)=0\qquad (i\notin d).
\]
If \(d=[0,x]\) we say that \(F\) has cohomological dimension \(\le x\); if \(d=[-x,0]\), homological dimension \(\le x\). The functor is of finite dimension if some finite interval works.
\end{statement}

\begin{statement}{Proposition 1.2.10}
Under the hypotheses of Proposition 1.2.7, assume in addition that \(RF:D^+(\A)\to D^+(\B)\), respectively \(LF:D^-(\A)\to D^-(\B)\), is of finite dimension. Then every unbounded complex of acyclic objects is split, and the corresponding derived functor exists on \(D(\A)\).
\end{statement}

This is Verdier's bounded-amplitude argument. Since cohomology of the image of a truncation can contribute only in a fixed finite band, the comparison between an unbounded complex and its truncations stabilizes degree by degree. Hence the functor on the unbounded derived category is obtained as the stable value of the functor on bounded truncations.

\section{Fibred categories in derived categories}

\subsection{Introduction}

The reader may skip this section on a first reading. Its purpose is to make precise the formalism underlying the ``trivial duality theorem'' and to remove the ambiguity, raised earlier by Artin, about the variance of derived functors under changes of ringed sites. We use the language of categories fibred over a base category, as in SGA 1, Expos\'e VI. A cleaved fibred category over a category \(C\) is equivalent to a pseudo-functor \(C^{\circ}\to\mathbf{Cat}\); the analogous statement for cofibrations corresponds to covariant pseudo-functors.

\begin{statement}{Definition 2.1.1}
If \(X\) is a ringed site, \(D(X)\) denotes the derived category of the abelian category of sheaves of left modules on \(X\). When the ring must be mentioned, we write \(D(X,A)\).
\end{statement}

If \(f:X\to Y\) is a morphism of ringed sites, then one has
\[
Rf_*:D^+(X)\to D^+(Y),
\qquad
Lf^*:D^-(Y)\to D^-(X).
\]
Passing to opposite categories, the first becomes a functor \(D^+(X)^{\circ}\to D^+(Y)^{\circ}\). Thus the collection of categories \(D^+(X)\), as \(X\) varies, is naturally both fibred and cofibrated in different senses, depending on whether one follows inverse image, direct image, or their opposites.

\begin{statement}{2.1.2}
Let \(C\) be the category of ringed sites under consideration. The assignment
\[
X\longmapsto D^+(X)
\]
together with \(Rf_*\) for morphisms \(f:X\to Y\), defines a pseudo-functor from \(C\) to triangulated categories, whenever the relevant boundedness hypotheses for right derived direct images are satisfied. The assignment \(X\mapsto D^-(X)\) with \(Lf^*\) defines a pseudo-functor from \(C^{\circ}\) to triangulated categories.
\end{statement}

The coherence isomorphisms are the canonical ones
\[
R(gf)_*\simeq Rg_*Rf_* ,
\qquad
L(fg)^*\simeq Lg^*Lf^*,
\]
obtained from the composition theorem for derived functors. The point of spelling this out is that later formulas, for instance projection and base change, are identities not in an isolated derived category but in these fibred systems.

\begin{statement}{2.1.3}
Given a cartesian square of ringed sites
\[
\begin{array}{ccc}
X'&\xrightarrow{g'}&X\\
\downarrow f'&&\downarrow f\\
Y'&\xrightarrow{g}&Y,
\end{array}
\]
there are canonical base-change morphisms
\[
Lg^*Rf_*K\longrightarrow Rf'_*Lg'^*K
\]
for \(K\in D^+(X)\) whenever both sides are defined. These morphisms are functorial in the square and compatible with composition of squares.
\end{statement}

They are built by deriving the adjunction map \(g^*f_*	o f'_*g'^*\). The main theorems of the preceding expos\'es identify hypotheses under which these morphisms are isomorphisms: proper base change for proper morphisms, smooth base change for smooth morphisms and prime-to-characteristic coefficients, and finite or affine cohomological-dimension bounds for the other cases.

\subsection{Fibred categories in triangulated categories}

\begin{statement}{Definition 2.2.1}
A category fibred in triangulated categories over \(C\) is a fibred category \(\mathcal D\to C\) such that every fibre \(\mathcal D_X\) is a triangulated category and every inverse-image functor attached to a cartesian arrow is exact. A cleaving is normalized if the inverse image attached to the identity of \(X\) is the identity functor and if the comparison isomorphisms for composition satisfy the usual associativity condition.
\end{statement}

Equivalently, a normalized cleaved fibred category in triangulated categories is a pseudo-functor \(C^{\circ}\to\mathbf{TriCat}\), where \(\mathbf{TriCat}\) denotes the 2-category of triangulated categories and exact functors. Dually, a cofibrated category in triangulated categories is a pseudo-functor \(C\to\mathbf{TriCat}\).

\begin{statement}{2.2.2}
Let \(\mathcal D\to C\) be fibred in triangulated categories. A morphism \(u:A	o B\) in the total category lying above \(f:X\to Y\) may be written, after choosing a cleaving, as a morphism
\[
A\longrightarrow f^*B
\]
in the fibre over \(X\). It is cartesian if this morphism is an isomorphism. Exactness of pullback means that cartesian liftings carry distinguished triangles to distinguished triangles.
\end{statement}

This language will be used to say that a construction is compatible with all base changes, rather than restating the corresponding square for each occurrence.

\begin{statement}{2.2.3}
Let \(\mathcal D\to C\) and \(\mathcal E\to C\) be fibred in triangulated categories. A morphism of fibred categories \(\Phi:\mathcal D\to\mathcal E\) is exact if its restriction to every fibre is exact and if it carries cartesian arrows to cartesian arrows. A natural transformation between such morphisms is compatible with base change if it is a morphism in the 2-category of fibred categories over \(C\).
\end{statement}

Thus the projection formula, the Kunneth formula, and the base-change formula will be recorded as morphisms, or isomorphisms, of functors between fibred triangulated categories.

\subsection{The trivial duality formula}

Let \(f:X\to Y\) be a morphism of ringed sites. Suppose for the moment that all objects involved are bounded in such a way that no ambiguity about derived functors occurs. There is a canonical morphism
\[
Lf^*R\mathcal Hom_Y(K,L)	o R\mathcal Hom_X(Lf^*K,Lf^*L).
\]
Adjunction gives the corresponding morphism
\[
R\mathcal Hom_Y(K,Rf_*M)	o Rf_*R\mathcal Hom_X(Lf^*K,M).
\]
This is called the trivial duality morphism: it is a formal consequence of the adjunction \((Lf^*,Rf_*)\), not of any properness or finiteness theorem.

\begin{statement}{Proposition 2.3.1}
Under the standard boundedness hypotheses, for \(K\in D^-(Y)\) and \(M\in D^+(X)\), there is a functorial morphism
\[
R\mathcal Hom_Y(K,Rf_*M)	o Rf_*R\mathcal Hom_X(Lf^*K,M).
\]
If \(K\) is represented by a bounded complex of flat modules, this morphism is computed by the ordinary sheaf Hom adjunction. It is compatible with composition of morphisms of ringed sites.
\end{statement}

\begin{statement}{Corollary 2.3.2}
If \(K\) is perfect, or more generally if it has finite Tor-amplitude and the usual boundedness conditions hold, the preceding morphism is an isomorphism. Equivalently, the derived projection formula
\[
K\otimes_Y^L Rf_*M\to Rf_*(Lf^*K\otimes_X^L M)
\]
is an isomorphism in the same range.
\end{statement}

The proof is local on \(Y\) and reduces to a finite complex of finite free modules, where it is the elementary tensor-Hom adjunction. The formulation in fibred categories guarantees compatibility with later gluing constructions.

\section{Gluing fibred categories and compactifiable morphisms}

\subsection{Introduction}

The construction of \(Rf_!\) for a non-proper morphism proceeds by choosing a compactification and setting \(Rf_!=R\bar f_*j_!\). To prove independence of the compactification, and to obtain functoriality, one needs a gluing formalism. This section records it in the language of fibred triangulated categories.

\subsection{Compactifiable morphisms}

\begin{statement}{Definition 3.2.1}
A morphism \(f:X\to S\) is compactifiable if it admits a factorization
\[
X\xrightarrow{j}\bar X\xrightarrow{\bar f}S
\]
where \(j\) is an open immersion and \(\bar f\) is proper. Such a factorization is called a compactification of \(f\).
\end{statement}

For schemes, separated morphisms of finite type are compactifiable in the noetherian situations considered here. The category of compactifications of a fixed compactifiable morphism is cofiltered after allowing domination: two compactifications are dominated by the closure of the diagonal image in their product over the base. This is the basic geometric fact behind independence of compactification.

\begin{statement}{Proposition 3.2.2}
Let \(f:X\to S\) be compactifiable. Given two compactifications \(X\hookrightarrow\bar X_1\to S\) and \(X\hookrightarrow\bar X_2\to S\), there is a third compactification \(X\hookrightarrow\bar X\to S\) and proper morphisms \(\bar X\to\bar X_i\) inducing the identity on \(X\). If the original compactifications are compatible with a base change, the dominating compactification may be chosen compatibly after refinement.
\end{statement}

The proof is the usual closure of the graph of \(X\) in \(\bar X_1\times_S\bar X_2\). Properness is stable under products and closed immersions, and the open immersion of \(X\) into the closure is induced by the diagonal.

\subsection{Gluing}

Let \(i:Z\hookrightarrow X\) be a closed immersion and \(j:U\hookrightarrow X\) the complementary open immersion. For abelian sheaves there is a short exact sequence
\[
0\to j_!j^*F\to F\to i_*i^*F\to0.
\]
In the derived category this becomes the localization triangle
\[
Rj_!j^*K\to K\to Ri_*i^*K\to Rj_!j^*K[1].
\]

\begin{statement}{Proposition 3.3.1}
Let \(X=U\cup Z\) be an open-closed decomposition. A functor \(T:D(X)\to\mathcal T\) to a triangulated category is determined, up to unique isomorphism, by its restrictions to \(D(U)\) and \(D(Z)\), together with compatibility with the localization triangle. More precisely, if two exact functors agree after composition with \(j_!\) and \(i_*\), and the comparison respects the boundary morphism, then they agree on all of \(D(X)\).
\end{statement}

Indeed every \(K\in D(X)\) sits in the localization triangle above. Applying the two functors and using the assumed agreement at the first and third vertices gives the agreement at the middle vertex by the triangulated five lemma; uniqueness follows from the same argument applied to morphisms.

\begin{statement}{Proposition 3.3.2}
Suppose \(f:X\to S\) is compactifiable and \(X=U\cup Z\) is an open-closed decomposition. If the functors \(Rf_!\) have been constructed for \(f|_U\) and \(f|_Z\), and if they satisfy base change and localization for those pieces, then there is a unique exact functor \(Rf_!:D(X)\to D(S)\) fitting into the triangle
\[
R(f|_U)_!j^*K\to Rf_!K\to R(f|_Z)_!i^*K\to.
\]
It is compatible with proper direct image when \(f\) is proper.
\end{statement}

This is the abstract gluing step. It will be applied inductively, using compactifications to replace open pieces by proper schemes and using the already established proper base-change theorem.

\begin{statement}{Proposition 3.3.3}
Let \(X\hookrightarrow\bar X\to S\) and \(X\hookrightarrow\bar X'\to S\) be two compactifications, with \(\bar X'\to\bar X\) proper and equal to the identity over \(X\). Then for every \(K\in D(X)\) there is a canonical isomorphism
\[
R\bar f_*j_!K\simeq R\bar f'_*j'_!K.
\]
These isomorphisms are transitive for further domination.
\end{statement}

Let \(p:\bar X'\to\bar X\) be the proper morphism. Since \(p\) is an isomorphism over \(X\), proper base change for the open immersion shows that
\[
j_!K\simeq Rp_*j'_!K.
\]
Applying \(R\bar f_*\) gives the result. Transitivity follows from the uniqueness of the adjunction morphisms and the associativity isomorphism \(R(\bar f p)_*\simeq R\bar f_*Rp_*\).

\section{Resolutions and application to the base-change arrow}

\subsection{Flat resolutions}

The construction of the base-change arrow for \(Rf_!\), and the proof of its independence from auxiliary choices, require resolutions which are well behaved with respect to inverse image, tensor product, and extension by zero. Flat resolutions provide this control for module-valued sheaves.

\begin{statement}{Definition 4.1.1}
Let \((X,A)\) be a ringed topos. A complex \(P\) of \(A\)-modules is called flat-adapted if each \(P^n\) is flat and if tensoring by \(P\) preserves quasi-isomorphisms in the range under consideration. A resolution \(P\to K\) is flat-adapted if \(P\) is flat-adapted and the morphism is a quasi-isomorphism.
\end{statement}

\begin{statement}{Proposition 4.1.2}
Every bounded-above complex of modules on a ringed topos admits a flat-adapted resolution. These resolutions may be chosen functorially enough to define \(-\otimes_A^L-\) and \(Lf^*\), and complexes of flat-adapted objects are split for these left derived functors.
\end{statement}

The proof uses the fact that every module is a quotient of a flat module, for example a sum of free modules. The class of flat modules satisfies the closure condition in Proposition 1.2.6: in a short exact sequence \(0\to P\to Q\to R\to0\) with \(Q\) and \(R\) flat and with the sequence pure in the required sense, \(P\) remains adapted for tensor computations. The general statement is obtained by replacing ``flat'' by the standard class of \(K\)-flat objects.

\begin{statement}{4.1.3}
If \(f:(X,A)\to(Y,B)\) is a morphism of ringed topoi, a flat-adapted resolution \(P\to K\) of a complex on \(Y\) computes
\[
Lf^*K=f^*P.
\]
Moreover, if \(Q\) is flat-adapted on \(X\), then \(f^*P\otimes_A Q\) computes \(Lf^*K\otimes_A^L Q\).
\end{statement}

These compatibilities are needed to make the projection morphism and the base-change morphism into morphisms of triangulated functors, rather than merely into maps after choosing arbitrary resolutions.

\subsection{Godement flasque resolutions}

Let \(X\) be a topos with enough points, or more generally work with a conservative family of fibre functors. The Godement construction associates to a sheaf \(F\) a cosimplicial flasque sheaf, and hence a resolution \(F\to C^\bullet(F)\). It is functorial in \(F\) and compatible with inverse image to the extent needed for base-change maps.

\begin{statement}{Definition 4.2.1}
A sheaf \(I\) on a topos \(X\) is flasque for the present purposes if the functor of sections over any object of the site has vanishing higher derived functors on \(I\). A complex is flasque-adapted if applying global sections, or more generally direct image, to it computes derived direct image.
\end{statement}

\begin{statement}{Proposition 4.2.2}
Every bounded-below complex of abelian sheaves admits a functorial flasque-adapted resolution. If \(f:X\to Y\) is a morphism of topoi, applying \(f_*\) to such a resolution computes \(Rf_*\). The construction is compatible with restriction to open subtopoi and with finite products of sites.
\end{statement}

For a single sheaf \(F\), embed it functorially into the product of its skyscraper sheaves at the chosen points; iterate the cokernel. The resulting resolution is flasque, and exactness is checked on the conservative family of points. For complexes, apply the construction by the usual Cartan-Eilenberg procedure. The functoriality is strict enough for all comparison morphisms used later.

\begin{statement}{4.2.3}
Let
\[
\begin{array}{ccc}
X'&\xrightarrow{g'}&X\\
\downarrow f'&&\downarrow f\\
Y'&\xrightarrow{g}&Y
\end{array}
\]
be a cartesian square of topoi. Choosing functorial flasque resolutions gives a natural base-change morphism
\[
g^*Rf_*K\to Rf'_*g'^*K.
\]
This construction agrees with the morphism obtained by deriving the underived adjunction map.
\end{statement}

The agreement follows because both constructions are induced by the same map on degree-zero sheaves and because the derived category is obtained by inverting quasi-isomorphisms. Functoriality with respect to pasting of squares follows from the functoriality of the Godement resolution.

\subsection{The base-change theorem}

\begin{statement}{Theorem 4.3.1}
For the morphisms and coefficients in the range of the preceding base-change theorems, the canonical morphism
\[
Lg^*Rf_*K\longrightarrow Rf'_*Lg'^*K
\]
constructed by the resolution formalism coincides with the base-change morphism studied in Expos\'es XII--XVI. In the proper case, for torsion coefficients, it is an isomorphism; in the smooth case, for torsion prime to the residual characteristics, it is an isomorphism; and in the finite cohomological-dimension range it is an isomorphism or monomorphism in the stated degrees.
\end{statement}

The proof is a comparison of constructions. The base-change morphism of Expos\'e XII was defined fibrewise and then glued by noetherian induction. The present morphism is defined by resolving and applying adjunction. On a strict localization of the base, the two maps are identified after applying the fibre functor, where both are the natural restriction map on etale cohomology. Since the relevant fibre functors are conservative, the two morphisms are the same. The assertions about being an isomorphism are exactly the theorems already proved.

\begin{statement}{Corollary 4.3.2}
Let \(f:X\to S\) be compactifiable, choose a compactification \(X\xrightarrow{j}\bar X\xrightarrow{\bar f}S\), and set
\[
Rf_!K=R\bar f_*j_!K.
\]
For every base change \(g:S'\to S\) there is a canonical morphism
\[
Lg^*Rf_!K\to Rf'_!Lg'^*K.
\]
It is independent of the chosen compactification. Under the proper-base-change hypotheses for \(\bar f\), it is an isomorphism in the same range.
\end{statement}

Independence follows from Proposition 3.3.3 and the transitivity of dominating compactifications. If a third compactification dominates two given ones, the two induced base-change morphisms agree after pullback to the dominating compactification; by proper base change for the domination morphism they therefore agree already on the original compactifications.

\begin{statement}{Corollary 4.3.3}
The functor \(Rf_!\), for compactifiable morphisms, is functorial in \(f\). If \(X\xrightarrow{f}Y\xrightarrow{g}Z\) are compactifiable, then there is a canonical isomorphism
\[
R(gf)_!\simeq Rg_!Rf_!.
\]
This isomorphism is compatible with base change and with the localization triangles for open-closed decompositions.
\end{statement}

Choose compactifications of \(f\) and \(g\), dominate the induced compactification of \(gf\), and apply the proper functoriality of ordinary direct image. The possible choices are compared through further domination. The localization compatibility is checked by reducing to the case of a proper morphism and an open-closed decomposition, where it is the usual exact triangle for \(j_!\) and \(i_*\).

\begin{remarktext}{Remark 4.3.4}
The theorem in this subsection is not yet the final construction of \(Rf_!\). It gives the base-change arrow and the independence of compactification in the bounded situations where the required resolutions and cohomological dimensions are available. The next section of Expos\'e XVII begins the full construction of direct image with proper supports and its comparison, finiteness, Kunneth, and symmetric-power properties.
\end{remarktext}


\section{Direct images with proper supports}

\subsection{The fundamental construction}

We now begin the construction indicated in Section 3. For a compactifiable morphism one wants a functor \(Rf_!\) which agrees with \(Rf_*\) when \(f\) is proper and with extension by zero when \(f\) is an open immersion. The only possible definition is obtained by decomposing a compactifiable morphism into these two elementary cases and then proving that the result is independent of the decomposition.

\begin{statement}{5.1.1}
Let \(X\) be a topos, let \(U\) be an open of \(X\), and let \(j:U\to X\) be the inclusion. The extension-by-the-empty-object functor for sheaves of sets on \(U\) is left adjoint to restriction \(j^*\). The same notation \(j_!\) is used for pointed sheaves, sheaves of groups, and modules over a ring sheaf. In the abelian and module cases this is extension by zero.
\end{statement}

If \(i:F\to X\) is the closed complement of \(U\), then \(j^*j_!\) is the identity and \(i^*j_!\) is the constant functor with value the initial, equivalently zero, object. These properties characterize \(j_!\). They imply exactness of \(j_!\) in the abelian cases, the natural monomorphism \(j_!\hookrightarrow j_*\), and the transitivity formula
\[
(j_1j_2)_!\simeq j_{1!}j_{2!}.
\]

\begin{statement}{Lemma 5.1.2}
Let \(f:X\to Y\) be a morphism of topoi, let \(V\subset Y\) be open, and let \(U=f^{-1}V\). For the square
\[
\begin{array}{ccc}
U&\xrightarrow{f'}&V\\
\downarrow j'&&\downarrow j\\
X&\xrightarrow{f}&Y,
\end{array}
\]
there is a unique base-change isomorphism
\[
f^*j_!\xrightarrow{\sim}j'_!f'^*
\]
compatible with the usual base-change morphism \(f^*j_*\to j'_*f'^*\).
\end{statement}

The adjoint map \(j^*f_*\to f'_*j'^*\) is plainly an isomorphism. Transposing it gives \(j'_!f'^*\simeq f^*j_!\). Compatibility with the inclusions into direct image is checked after restricting to \(U\), where it is tautological.

\begin{statement}{5.1.3}
If \(X\) is ringed by \(A\), exactness of \(j_!\) gives an exact functor
\[
j_!:D(U,j^*A)\to D(X,A).
\]
The transitivity formula for open immersions and the adjunction formula
\[
\Hom_{D(X)}(j_!K,L)\simeq \Hom_{D(U)}(K,j^*L)
\]
hold in the derived categories.
\end{statement}

\begin{statement}{5.1.4}
Let \(f:X\to Y\) be proper, let \(A\) be a torsion ring sheaf on \(Y\), and ring \(X\) by \(f^*A\). If the dimensions of the fibres of \(f\) are bounded by \(n\), then \(R^qf_*F=0\) for \(q>2n\) and every \(f^*A\)-module \(F\). Hence
\[
Rf_*:D(X,f^*A)\to D(Y,A)
\]
is defined on the unbounded derived category in the range used here. For a composite of proper morphisms, one has the transitivity isomorphism \(R(gh)_*\simeq Rg_*Rh_*\).
\end{statement}

\begin{statement}{5.1.5--5.1.6}
Consider a commutative square
\[
\begin{array}{ccc}
U&\xrightarrow{j_1}&X\\
\downarrow g&&\downarrow f\\
V&\xrightarrow{j_2}&Y,
\end{array}
\]
where \(j_1,j_2\) are open immersions and \(f,g\) are proper of bounded fibre dimension. For \(K\in D(U)\) there is a canonical isomorphism
\[
j_{2!}Rg_*K\xrightarrow{\sim}Rf_*j_{1!}K.
\]
\end{statement}

After restricting to \(V\), the assertion becomes the proper base-change isomorphism for the cartesian square \(V\times_YX\). On the closed complement \(Y-V\), both sides vanish: for the right side this follows from proper base change and the fact that the restriction of \(j_{1!}K\) to the complementary closed fibre is zero. Thus the morphism is an isomorphism globally.

\begin{statement}{Theorem 5.1.8}
Let \(S\) be ringed by a torsion ring sheaf \(A\), and let \((S)\) be the category of quasi-compact quasi-separated \(S\)-schemes with \(S\)-compactifiable morphisms. There is, in one and essentially only one way, a functor
\[
Rf_!:D(X,A_X)\to D(Y,A_Y)
\]
for every morphism \(f:X\to Y\) in \((S)\), together with transitivity isomorphisms
\[
R(gf)_!\simeq Rg_!Rf_!,
\]
satisfying the cocycle condition. On proper morphisms it is \(Rf_*\); on open immersions it is \(j_!\); and for a square as in 5.1.5 the compatibility is the isomorphism just described.
\end{statement}

The existence is exactly the gluing theorem of Section 3 applied to the two subcategories of proper morphisms and open immersions. The uniqueness follows from the fact that every compactifiable morphism is, by definition, a composite of an open immersion and a proper morphism, and any two compactifications are dominated by a third. The compatibility for a dominated compactification is supplied by proper base change.

\begin{remarktext}{5.1.9}
When \(f:X\to\Spec k\), with \(k\) separably closed, one writes
\[
R^qf_!K=H^q_c(X,K)
\]
and calls this hypercohomology, or cohomology, with proper support. The notation \(Rf_!\) is slightly abusive: it is not the right derived functor of an ordinary functor \(f_!\) in general, although an ordinary \(f_!\) for pointed sheaves and groups will be described later.
\end{remarktext}

\begin{statement}{Variants 5.1.11--5.1.15}
The same construction gives: an ordinary functor \(f_!\) for pointed sheaves; a functor \(R^1f_!\) for sheaves of ind-finite groups, independent of compactification; a triangulated way-out functor
\[
Rf_!:D^+_{\mathrm{tors}}(X)\to D^+_{\mathrm{tors}}(Y);
\]
and the corresponding version for ringed schemes with torsion coefficient rings. These variants use only proper base change and the equality \(j_!=j_*\) for an open immersion which is also proper.
\end{statement}

\begin{statement}{5.1.16}
Let \(j:U\to X\) be open, let \(i:F\to X\) be the closed complement, and let \(f:X\to S\) be compactifiable. For \(K\in D(X,A_X)\), the localization triangle gives
\[
R(fj)_!j^*K\to Rf_!K\to R(fi)_!i^*K\to R(fj)_!j^*K[1].
\]
Consequently there is a long exact cohomology sequence with proper supports. Over a separably closed field this is the usual exact sequence
\[
\cdots\to H^q_c(U,G)\to H^q_c(X,G)\to H^q_c(F,G)\to H^{q+1}_c(U,G)\to\cdots .
\]
\end{statement}

\begin{statement}{5.1.17}
For a proper morphism \(u:X\to Y\) over \(S\), and \(K\in D(Y,A_Y)\), the adjunction map \(K\to Ru_*u^*K\) gives a contravariant pullback on cohomology with proper supports:
\[
u^*:Rg_!K\to Rf_!u^*K,
\]
where \(f=g u\). For \(S=\Spec k\), this is the map
\[
H^q_c(Y,K)\to H^q_c(X,u^*K).
\]
\end{statement}


\subsection{The base-change theorem}

\begin{statement}{5.2.1}
Let \(f:X\to Y\) be a morphism of topoi, let \(V\subset Y\) be open, and let \(U=f^{-1}V\). Write \(j:V\hookrightarrow Y\), \(j':U\hookrightarrow X\), and \(f':U\to V\). Let \(A\) be a ring sheaf on \(Y\), let \(A'\) be a ring sheaf on \(X\), and let \(L\) be a bounded-above complex of \((A',f^*A)\)-bimodules. There is a canonical isomorphism
\[
L\otimes^L_{f^*A} f^*j_!K
 \xrightarrow{\sim}
 j'_!\bigl(j'^*L\otimes^L_{f'^*A}f'^*K\bigr)
\]
for \(K\in D^-(V,A)\). If \(L\) has finite Tor-amplitude over \(f^*A\), the same formula is defined for arbitrary \(K\in D(V,A)\).
\end{statement}

The first ingredient is the base-change isomorphism \(f^*j_!\simeq j'_!f'^*\) of Lemma 5.1.2. The second is the elementary compatibility of extension by zero with tensor products:
\[
L\otimes^L_{f^*A}j'_!M
 \simeq
j'_!\bigl(j'^*L\otimes^L_{f'^*A}M\bigr).
\]
For complexes adapted to the relevant tensor functors this is an underived identity; the displayed formula follows by passage to derived categories.

\begin{statement}{5.2.2}
Let \(f:X\to S\) be a compactifiable morphism and choose a compactification \(X\xrightarrow{j}\bar X\xrightarrow{p}S\). For a base change \(g:S'\to S\), form the cartesian diagram
\[
\begin{array}{ccccc}
X'&\xrightarrow{g'}&X&\xrightarrow{j}&\bar X\\
\downarrow f'&&\downarrow f&&\downarrow p\\
S'&\xrightarrow{g}&S&=&S,
\end{array}
\]
with \(X'\xrightarrow{j'}\bar X'\xrightarrow{p'}S'\) obtained by base change. If \(A\) and \(A'\) are torsion coefficient rings on \(S\) and \(S'\), and if \(L\) is a bounded-above complex of \((A',g^*A)\)-bimodules, then the proper base-change theorem for \(p\) and the formula of 5.2.1 give an isomorphism
\[
L\otimes^L_{g^*A}g^*Rp_*j_!K
\;\xrightarrow{\sim}\;
Rp'_*j'_!\bigl(f'^*L\otimes^L_{g'^*A}g'^*K\bigr).
\]
Equivalently,
\[
L\otimes^L g^*Rf_!K\xrightarrow{\sim}Rf'_!\bigl(f'^*L\otimes^Lg'^*K\bigr)
\]
after identifying \(Rf_!=Rp_*j_!\) and \(Rf'_!=Rp'_*j'_!\).
\end{statement}

This is the base-change isomorphism for cohomology with proper supports attached to the chosen compactification. The construction is simply the composition of two already constructed isomorphisms: first move \(g^*\) past the proper direct image \(Rp_*\), and then move extension by zero past the pullback and tensor product.

\begin{statement}{Lemma 5.2.3}
The isomorphism of 5.2.2 is independent of the compactification \(X\hookrightarrow\bar X\to S\).
\end{statement}

It is enough to compare two compactifications when one dominates the other. Suppose \(q:\bar X_1\to\bar X_2\) is proper and induces the identity on \(X\). The isomorphism attached to \(\bar X_1\) is the composite of those attached to the open immersion \(j_1\), the proper map \(q\), and the proper map \(p_2\); the isomorphism attached to \(\bar X_2\) is the composite of those attached to \(j_2\) and \(p_2\). By adjunction for \(j_{2!}\), equality can be checked after restriction to the open \(X\), where both composites are the identity. Two arbitrary compactifications are dominated by a third, so the assertion follows.

\begin{statement}{Lemma 5.2.4}
Let \(X\xrightarrow{v}Y\xrightarrow{u}S\) be a composite of compactifiable morphisms. The base-change isomorphisms of 5.2.2 for \(u\), \(v\), and \(uv\) are compatible with the transitivity isomorphism
\[
R(uv)_!\simeq Ru_!Rv_!.
\]
In other words, the two natural ways of passing from
\[
L\otimes^L g^*R(uv)_!K
\]
to the corresponding expression on the base-changed composite agree.
\end{statement}

Choose compatible compactifications of \(v\) and \(u\), and hence of their composite. The assertion then reduces to the associativity of proper base change and to the transitivity formula for extension by zero. The general case follows by replacing the chosen compactifications by a common dominating system and using Lemma 5.2.3.



\begin{remarktext}{Further compatibility 5.2.4.1}
The pullback just described is not an isolated operation. If \(X\xrightarrow{u}Y\xrightarrow{v}Z\) are proper over \(S\), the pullback for \(vu\) is the composite of the pullbacks for \(v\) and \(u\). If an open immersion is inserted in the diagram, the same assertion remains true after replacing ordinary direct image by extension by zero. These compatibilities are immediate on a chosen compactification and therefore hold in the compactification-independent formalism by the uniqueness statement of Theorem 5.1.8.
\end{remarktext}

\begin{statement}{5.2.4.2}
Let \(K\to L\to M\to K[1]\) be a distinguished triangle on \(X\). For a compactifiable \(f:X\to S\), the triangle
\[
Rf_!K\to Rf_!L\to Rf_!M\to Rf_!K[1]
\]
commutes with all pullbacks by proper morphisms over \(S\) and with the localization triangle for an open-closed decomposition of \(X\). Thus the exact sequences with proper support are natural not only in the coefficient complex but also in the geometric pair.
\end{statement}

\begin{remarktext}{Transition within 5.2}
The construction so far gives the functor and its formal exactness. What is still missing is the strong base-change theorem: under the hypotheses used in \(\ell\)-adic and torsion cohomology, the base-change morphism
\[
g^*Rf_!K\to Rf'_!g'^*K
\]
should be an isomorphism. Section 5.2 is devoted to this statement and to its reduction to the proper and open cases already isolated above.
\end{remarktext}

\bigskip


\end{document}
